\chapter{ระบบอะตอมในสามมิติ โมเมนตัมเชิงมุม และสปินทรอนิกส์}
\label{chap:hydrogen_atom_spin}
\index{อะตอมไฮโดรเจน (Hydrogen Atom)}
\index{โมเมนตัมเชิงมุมออร์บิทัล (Orbital Angular Momentum)}
\index{ฮาร์มอนิกส์ทรงกลม (Spherical Harmonics)}
\index{การทดลองชแตร์น-แกร์ลัค (Stern-Gerlach Experiment)}
\index{เมทริกซ์สปินเพาลี (Pauli Spin Matrices)}
\index{สปินทรอนิกส์ (Spintronics)}
\index{โครงสร้างละเอียด (Fine Structure)}
\index{การสร้างภาพด้วยเรโซแนนซ์แม่เหล็ก (MRI)}

% ==========================================================
% แผนบริหารการสอนประจำบทที่ 6
% ==========================================================
\section*{แผนบริหารการสอนประจำบทที่ 6}
\addcontentsline{toc}{section}{แผนบริหารการสอนประจำบทที่ 6}

\begin{lessonplanbox}
\noindent\textbf{หัวข้อเนื้อหาประจำบท}
\begin{enumerate}
    \item สมการชเรอดิงเงอร์ในสามมิติและพิกัดทรงกลม
    \item ตัวดำเนินการโมเมนตัมเชิงมุม พีชคณิตคอมมิวเตเตอร์ และฮาร์มอนิกส์ทรงกลม
    \item อะตอมไฮโดรเจนและการแยกตัวแปรในแนวรัศมี
    \item ศักย์ยังผล การกระจายอนุกรมโฟรเบนิอุส และการควอนไตซ์พลังงาน
    \item ฟังก์ชันคลื่น รูปร่างของออร์บิทัล บัพ และค่าคาดหวังเชิงรัศมี
    \item การทดลองชแตร์น-แกร์ลัคและหลักฐานเชิงประจักษ์ของสปินอิเล็กตรอน
    \item พีชคณิตสปิน เมทริกซ์เพาลี และสปินอร์
    \item การหมุนของสปินอร์และความสมมาตรคาบ $4\pi$
    \item อันตรกิริยาสปิน-ออร์บิท โครงสร้างละเอียด และปรากฏการณ์ซีมาน
    \item การประยุกต์เทคโนโลยีควอนตัมสมัยใหม่ ---- สปินทรอนิกส์ วัสดุ 2 มิติ และ MRI
    \item ห้องปฏิบัติการเสมือนจริง 3 มิติ WebAR สำรวจออร์บิทัลและชแตร์น-แกร์ลัค
    \item ตัวอย่างการแก้โจทย์เชิงวิเคราะห์ตามระเบียบวิธี P-S-C-T
    \item สคริปต์การคำนวณเชิงตัวเลขภาษา Python จำลองออร์บิทัลสามมิติ
    \item สรุปสาระสำคัญและแบบฝึกหัดท้ายบท
\end{enumerate}

\noindent\textbf{วัตถุประสงค์เชิงพฤติกรรม}
\begin{enumerate}
    \item แปลงตัวดำเนินการลาปลาเซียนและแยกตัวแปรสมการชเรอดิงเงอร์ในพิกัดทรงกลมสำหรับศักย์สู่ศูนย์กลางได้
    \item พิสูจน์ความสัมพันธ์คอมมิวเตเตอร์ของโมเมนตัมเชิงมุมและหาค่าเจาะจงของ $\hat{\mathbf{L}}^2$ และ $\hat{L}_z$ ด้วยฮาร์มอนิกส์ทรงกลมได้
    \item แก้สมการรัศมีของอะตอมไฮโดรเจนด้วยระเบียบวิธีอนุกรมและอธิบายที่มาของเลขควอนตัมหลัก $n$ ได้
    \item คำนวณความหนาแน่นความน่าจะเป็นในแนวรัศมี จำนวนบัพ และค่าคาดหวัง $\langle r \rangle, \langle r^2 \rangle, \langle 1/r \rangle$ ได้
    \item อธิบายผลการทดลองชแตร์น-แกร์ลัคและวิเคราะห์การวัดสปินแบบเรียงลำดับได้
    \item ประยุกต์ใช้เมทริกซ์สปินเพาลี สปินอร์ และตัวดำเนินการหมุนสปินในการคำนวณความน่าจะเป็นของสถานะควอนตัมได้
    \item วิเคราะห์การเลื่อนของระดับพลังงานจากอันตรกิริยาสปิน-ออร์บิทและโครงสร้างละเอียดได้
    \item อธิบายหลักการทำงานของเทคโนโลยีสปินทรอนิกส์ GMR, MRAM และเครื่องตรวจวินิจฉัยทางการแพทย์ MRI ได้
\end{enumerate}
\end{lessonplanbox}

% ==========================================================
% 6.1 สมการชเรอดิงเงอร์ในสามมิติและพิกัดทรงกลม
% ==========================================================
\newpage
\section{สมการชเรอดิงเงอร์ในสามมิติและพิกัดทรงกลม}
\index{พิกัดทรงกลม (Spherical Coordinates)}
\index{ศักย์สู่ศูนย์กลาง (Central Potential)}

ในโลกกายภาพความเป็นจริง ปรากฏการณ์ควอนตัมส่วนใหญ่ดำเนินไปในปริภูมิสามมิติ ($\mathbb{R}^3$) ไม่ว่าจะเป็นการโคจรของอิเล็กตรอนรอบนิวเคลียส การกระเจิงของอนุภาคในสนามแรง หรือพันธะเคมีระหว่างอะตอม สมการชเรอดิงเงอร์ไม่ขึ้นกับเวลาในระบบพิกัดคาร์ทีเซียน $(x, y, z)$ เขียนได้ในรูปสากลดังนี้
\begin{equation}
-\frac{\hbar^2}{2m} \nabla^2 \psi(\mathbf{r}) + V(\mathbf{r}) \psi(\mathbf{r}) = E \psi(\mathbf{r})
\label{eq:tise_3d_cartesian}
\end{equation}
โดยที่ตัวดำเนินการลาปลาเซียน (Laplacian Operator) ในพิกัดคาร์ทีเซียนมีรูปทรงผลรวมอนุพันธ์อันดับสอง
\begin{equation}
\nabla^2 = \frac{\partial^2}{\partial x^2} + \frac{\partial^2}{\partial y^2} + \frac{\partial^2}{\partial z^2}
\end{equation}

ทว่าสำหรับระบบกายภาพที่มีความสำคัญพื้นฐานที่สุดในธรรมชาติ เช่น แรงดึงดูดทางไฟฟ้าสถิตคูลอมบ์ระหว่างโปรตอนและอิเล็กตรอนในอะตอมไฮโดรเจน พลังงานศักย์ขึ้นอยู่กับระยะห่างสัมพัทธ์ระหว่างอนุภาคเพียงอย่างเดียว โดยไม่ขึ้นกับมุมทิศทาง กล่าวคือ $V(\mathbf{r}) = V(r)$ ศักย์ลักษณะนี้เรียกว่า \textbf{ศักย์สู่ศูนย์กลาง (Central Potential)} ซึ่งมีสมมาตรการหมุนทรงกลมอย่างสมบูรณ์รอบจุดกำเนิด (Rotational Invariance) การแก้ปัญหาในพิกัดคาร์ทีเซียนจะนำไปสู่สมการอนุพันธ์ย่อยที่ตัวแปร $x, y, z$ พันกันอย่างซับซ้อน ดังนั้น การเปลี่ยนไปใช้ \textbf{ระบบพิกัดทรงกลม (Spherical Polar Coordinates, $(r, \theta, \phi)$)} จึงเป็นระเบียบวิธีที่ธรรมชาติของระบบเรียกร้อง

\begin{figure}[H]
\centering
\begin{tikzpicture}[scale=1.1, >=Stealth]
    % Coordinate axes
    \draw[->, thick, slateGrey] (0, 0, 0) -- (4.0, 0, 0) node[right, font=\small\bfseries] {$y$};
    \draw[->, thick, slateGrey] (0, 0, 0) -- (0, 3.8, 0) node[above, font=\small\bfseries] {$z$};
    \draw[->, thick, slateGrey] (0, 0, 0) -- (0, 0, 4.0) node[below left, font=\small\bfseries] {$x$};

    % Point P coordinates
    \coordinate (O) at (0, 0, 0);
    \coordinate (P) at (2.2, 2.8, 2.0);
    \coordinate (Pproj) at (2.2, 0, 2.0);

    % Projection lines
    \draw[dashed, slateGrey!80] (O) -- (Pproj);
    \draw[dashed, slateGrey!80] (P) -- (Pproj);
    \draw[dashed, slateGrey!80] (Pproj) -- (0, 0, 2.0);
    \draw[dashed, slateGrey!80] (Pproj) -- (2.2, 0, 0);

    % Vector r
    \draw[->, line width=1.5pt, cyberBlue] (O) -- (P) node[midway, above left, font=\normalsize\bfseries] {$\mathbf{r}$};
    \filldraw[rubyRed] (P) circle (2.2pt) node[above right, font=\small\bfseries] {$P(r, \theta, \phi)$};

    % Angles
    % Theta (polar angle from z-axis)
    \draw[thick, entanglePink, ->] (0, 1.2, 0) arc (90:40:1.2);
    \node at (0.35, 1.5, 0.4) [font=\small\bfseries\color{entanglePink}] {$\theta$};

    % Phi (azimuthal angle from x-axis in xy plane)
    \draw[thick, quantumGold, ->] (0, 0, 1.5) arc (0:45:1.5);
    \node at (0.9, 0, 1.3) [font=\small\bfseries\color{quantumGold}] {$\phi$};
\end{tikzpicture}
\caption{ระบบพิกัดทรงกลมแสดงเวกเตอร์ตำแหน่ง $\mathbf{r}$ ระยะรัศมี $r$ มุมขั้ว $\theta$ และมุมทิศ $\phi$}
\label{fig:spherical_coordinates}
\end{figure}

ความสัมพันธ์ระหว่างพิกัดคาร์ทีเซียนและพิกัดทรงกลมกำหนดโดย
\begin{align}
x &= r \sin\theta \cos\phi \label{eq:cart_to_sph_x} \\
y &= r \sin\theta \sin\phi \label{eq:cart_to_sph_y} \\
z &= r \cos\theta \label{eq:cart_to_sph_z}
\end{align}
โดยมีขอบเขตของตัวแปรในปริภูมิสามมิติดังนี้
\begin{equation}
r \in [0, \infty), \qquad \theta \in [0, \pi], \qquad \phi \in [0, 2\pi)
\end{equation}

\subsection{การแปลงตัวดำเนินการลาปลาเซียนในพิกัดทรงกลม}
ในการแปลงตัวดำเนินการลาปลาเซียน $\nabla^2$ ให้อยู่ในพิกัดเชิงเส้นโค้งตั้งฉาก (Orthogonal Curvilinear Coordinates) เราพิจารณาองค์ประกอบส่วนเพิ่มของความยาวระยะทางกำลังสอง (Line Element Squared)
\begin{equation}
ds^2 = dx^2 + dy^2 + dz^2 = dr^2 + r^2 d\theta^2 + r^2 \sin^2\theta d\phi^2
\end{equation}
ทำให้ได้ตัวประกอบมาตราส่วน (Scale Factors) สำหรับแต่ละแกนพิกัดดังนี้
\begin{equation}
h_r = 1, \qquad h_\theta = r, \qquad h_\phi = r\sin\theta
\end{equation}

จากสูตรทั่วไปของลาปลาเซียนในระบบพิกัดตั้งฉากที่มีตัวประกอบมาตราส่วน $h_1, h_2, h_3$
\begin{equation}
\nabla^2 \psi = \frac{1}{h_1 h_2 h_3} \sum_{i=1}^3 \frac{\partial}{\partial u_i} \left( \frac{h_1 h_2 h_3}{h_i^2} \frac{\partial \psi}{\partial u_i} \right)
\end{equation}
เมื่อแทนค่า $u_1 = r, u_2 = \theta, u_3 = \phi$ และ $h_r h_\theta h_\phi = r^2 \sin\theta$ จะได้
\begin{equation}
\nabla^2 = \frac{1}{r^2} \frac{\partial}{\partial r}\left( r^2 \frac{\partial}{\partial r} \right) + \frac{1}{r^2 \sin\theta} \frac{\partial}{\partial \theta}\left( \sin\theta \frac{\partial}{\partial \theta} \right) + \frac{1}{r^2 \sin^2\theta} \frac{\partial^2}{\partial \phi^2}
\label{eq:laplacian_spherical_full}
\end{equation}

อนุพันธ์ย่อยในสมการ~\eqref{eq:laplacian_spherical_full} สามารถจัดกลุ่มอย่างเป็นระบบออกเป็นสองส่วนอย่างชัดเจน ส่วนแรกคือ \textbf{ส่วนอนุพันธ์เชิงรัศมี (Radial Part)} และส่วนหลังคือ \textbf{ส่วนอนุพันธ์เชิงมุม (Angular Part)}
\begin{equation}
\nabla^2 = \frac{1}{r^2} \frac{\partial}{\partial r}\left( r^2 \frac{\partial}{\partial r} \right) + \frac{1}{r^2} \nabla_{\theta, \phi}^2
\end{equation}
โดยที่ตัวดำเนินการลาปลาเซียนเชิงมุมนิยามโดย
\begin{equation}
\nabla_{\theta, \phi}^2 \equiv \frac{1}{\sin\theta} \frac{\partial}{\partial \theta}\left( \sin\theta \frac{\partial}{\partial \theta} \right) + \frac{1}{\sin^2\theta} \frac{\partial^2}{\partial \phi^2}
\label{eq:angular_laplacian}
\end{equation}

\subsection{การแยกตัวแปรเชิงรัศมีและเชิงมุม}
เราตั้งสมมติฐานว่า ฟังก์ชันคลื่นรวมสามารถแยกตัวแปรแบบผลคูณได้ (Separation of Variables)
\begin{equation}
\psi(r, \theta, \phi) = R(r) Y(\theta, \phi)
\label{eq:psi_separation_ry}
\end{equation}
เมื่อแทนสมการ~\eqref{eq:psi_separation_ry} ลงในสมการชเรอดิงเงอร์สามมิติ~\eqref{eq:tise_3d_cartesian} ที่ใช้ลาปลาเซียนพิกัดทรงกลม
\begin{equation}
-\frac{\hbar^2}{2m} \left[ \frac{Y}{r^2} \frac{d}{dr}\left( r^2 \frac{dR}{dr} \right) + \frac{R}{r^2} \nabla_{\theta, \phi}^2 Y \right] + V(r) R Y = E R Y
\end{equation}

คูณตลอดทั้งสมการด้วย $-\frac{2mr^2}{\hbar^2}$ และหารด้วย $R(r)Y(\theta, \phi)$ จะได้สมการที่ตัวแปรแยกออกจากกันอย่างสมบูรณ์
\begin{equation}
\frac{1}{R(r)} \frac{d}{dr}\left( r^2 \frac{dR}{dr} \right) - \frac{2m r^2}{\hbar^2}\big[ V(r) - E \big] = -\frac{1}{Y(\theta, \phi)} \nabla_{\theta, \phi}^2 Y(\theta, \phi)
\label{eq:separated_master}
\end{equation}

พิจารณาสมการ~\eqref{eq:separated_master} จะพบว่า ด้านซ้ายมือของสมการเป็นฟังก์ชันของตัวแปร $r$ เพียงตัวเดียว ในขณะที่ด้านขวามือเป็นฟังก์ชันของตัวแปรงมุม $(\theta, \phi)$ เพียงอย่างเดียว การที่สองด้านนี้จะเท่ากันได้ในทุกๆ จุดของปริภูมิสามมิติ ทั้งสองด้านจะต้องเท่ากับค่าคงตัวเดียวกัน ค่าคงตัวนี้เรียกว่า \textbf{ค่าคงตัวของการแยกตัวแปร (Separation Constant)} ซึ่งเพื่อความสอดคล้องกับคุณสมบัติของโมเมนตัมเชิงมุม เรากำหนดให้ค่าคงตัวนี้มีค่าเท่ากับ $\ell(\ell + 1)$ โดยที่ $\ell$ เป็นจำนวนเต็มไม่เป็นลบ

ส่งผลให้สมการแยกออกเป็นสองสมการเอกเทศ
\begin{align}
\nabla_{\theta, \phi}^2 Y(\theta, \phi) &= -\ell(\ell + 1) Y(\theta, \phi) \label{eq:angular_eigen_eq} \\
\frac{d}{dr}\left( r^2 \frac{dR}{dr} \right) - \frac{2m r^2}{\hbar^2}\big[ V(r) - E \big] R(r) &= \ell(\ell + 1) R(r) \label{eq:radial_general_eq}
\end{align}

ความจริงอันน่าทึ่งยิ่งของกลศาสตร์ควอนตัมคือ สมการเชิงมุม~\eqref{eq:angular_eigen_eq} มีรูปแบบที่เหมือนกันสำหรับทุกๆ ศักย์สู่ศูนย์กลาง $V(r)$ ไม่ว่าจะเป็นศักย์คูลอมบ์ของอะตอมไฮโดรเจน ศักย์ยูกาวาของแรงนิวเคลียร์ หรือศักย์ตัวแกว่งกวัดฮาร์มอนิก 3 มิติ ผลเฉลยเชิงมุมจะขึ้นกับความสมมาตรทรงกลมและโมเมนตัมเชิงมุมเสมอ

% ==========================================================
% 6.2 ตัวดำเนินการโมเมนตัมเชิงมุมและฮาร์มอนิกส์ทรงกลม
% ==========================================================
\section{ตัวดำเนินการโมเมนตัมเชิงมุมและฮาร์มอนิกส์ทรงกลม}
\index{โมเมนตัมเชิงมุมออร์บิทัล (Orbital Angular Momentum)}
\index{ฮาร์มอนิกส์ทรงกลม (Spherical Harmonics)}
\index{ตัวดำเนินการเพิ่มลด (Ladder Operators)}

ในกลศาสตร์คลาสสิก โมเมนตัมเชิงมุมนิยามโดยผลคูณเชิงเวกเตอร์ $\mathbf{L} = \mathbf{r} \times \mathbf{p}$ ในกลศาสตร์ควอนตัม เราแทนที่เวกเตอร์ตำแหน่งและโมเมนตัมด้วยตัวดำเนินการเชิงเส้นบนปริภูมิฮิลเบิร์ต $\hat{\mathbf{r}}$ และ $\hat{\mathbf{p}} = -i\hbar\boldsymbol{\nabla}$ ทำให้ได้ \textbf{ตัวดำเนินการโมเมนตัมเชิงมุมออร์บิทัล (Orbital Angular Momentum Operator)}
\begin{equation}
\hat{\mathbf{L}} = \hat{\mathbf{r}} \times \hat{\mathbf{p}}
\end{equation}
ซึ่งมีองค์ประกอบตามแนวแกนพิกัดคาร์ทีเซียนดังนี้
\begin{align}
\hat{L}_x &= \hat{y}\hat{p}_z - \hat{z}\hat{p}_y = -i\hbar\left( y\frac{\partial}{\partial z} - z\frac{\partial}{\partial y} \right) \\
\hat{L}_y &= \hat{z}\hat{p}_x - \hat{x}\hat{p}_z = -i\hbar\left( z\frac{\partial}{\partial x} - x\frac{\partial}{\partial z} \right) \\
\hat{L}_z &= \hat{x}\hat{p}_y - \hat{y}\hat{p}_x = -i\hbar\left( x\frac{\partial}{\partial y} - y\frac{\partial}{\partial x} \right)
\end{align}

\subsection{พีชคณิตคอมมิวเตเตอร์ของโมเมนตัมเชิงมุม}
การสลับที่ระหว่างตัวดำเนินการโมเมนตัมเชิงมุมแสดงถึงความสัมพันธ์ทางเรขาคณิตที่ลึกซึ้งของการหมุนในสามมิติ โดยใช้ความสัมพันธ์คอมมิวเตเตอร์คาโนนิคัล $[\hat{x}_i, \hat{p}_j] = i\hbar\delta_{ij}$
\begin{align*}
[\hat{L}_x, \hat{L}_y] &= [\hat{y}\hat{p}_z - \hat{z}\hat{p}_y, \, \hat{z}\hat{p}_x - \hat{x}\hat{p}_z] \\
&= [\hat{y}\hat{p}_z, \hat{z}\hat{p}_x] + [\hat{z}\hat{p}_y, \hat{x}\hat{p}_z] \\
&= \hat{y}\hat{p}_x [\hat{p}_z, \hat{z}] + \hat{p}_y\hat{x} [\hat{z}, \hat{p}_z] \\
&= \hat{y}\hat{p}_x (-i\hbar) + \hat{x}\hat{p}_y (i\hbar) = i\hbar(\hat{x}\hat{p}_y - \hat{y}\hat{p}_x) = i\hbar \hat{L}_z
\end{align*}

โดยการสลับดัชนีแบบวนรอบ (Cyclic Permutation) เราสรุปพีชคณิตลี $\mathfrak{so}(3)$ ของโมเมนตัมเชิงมุมได้ดังนี้
\begin{equation}
[\hat{L}_i, \hat{L}_j] = i\hbar \sum_{k=1}^3 \epsilon_{ijk} \hat{L}_k
\label{eq:angmom_lie_algebra}
\end{equation}
โดยที่ $\epsilon_{ijk}$ คือสัญลักษณ์เลวี-ชีวิตา (Levi-Civita Symbol)

เนื่องจากองค์ประกอบแต่ละแกนไม่สลับที่กันเอง ($[\hat{L}_x, \hat{L}_y] \neq 0$) ตามหลักความไม่แน่นอนของไฮเซนแบร์ก เราจึงไม่สามารถวัดค่าโมเมนตัมเชิงมุมตามแกนที่แตกต่างกันได้พร้อมกันอย่างแม่นยำแน่นอน ทว่าเมื่อพิจารณา \textbf{ขนาดกำลังสองของโมเมนตัมเชิงมุมรวม}
\begin{equation}
\hat{\mathbf{L}}^2 = \hat{L}_x^2 + \hat{L}_y^2 + \hat{L}_z^2
\end{equation}
เราสามารถพิสูจน์ได้โดยตรงว่า $\hat{\mathbf{L}}^2$ สลับที่กับองค์ประกอบทุกแกน
\begin{equation}
[\hat{\mathbf{L}}^2, \hat{L}_x] = [\hat{\mathbf{L}}^2, \hat{L}_y] = [\hat{\mathbf{L}}^2, \hat{L}_z] = 0
\label{eq:l2_commute_all}
\end{equation}

ความสัมพันธ์นี้มีความหมายทางฟิสิกส์อันลึกซึ้งยิ่ง กล่าวคือ เราสามารถค้นหาสถานะควอนตัมร่วมที่ทำให้ขนาดกำลังสองรวม $\hat{\mathbf{L}}^2$ และองค์ประกอบตามแนวแกนใดแกนหนึ่ง (ธรรมเนียมนิยมสากลเลือกแกน $z$ คือ $\hat{L}_z$) มีค่าแน่นอนพร้อมกันได้อย่างแม่นตรง

\subsection{ตัวดำเนินการในพิกัดทรงกลมและฮาร์มอนิกส์ทรงกลม}
เมื่อแปลงอนุพันธ์คาร์ทีเซียนเป็นพิกัดทรงกลม ตัวดำเนินการ $\hat{L}_z$ และ $\hat{\mathbf{L}}^2$ จะมีรูปที่เรียบง่ายอย่างน่าอัศจรรย์
\begin{align}
\hat{L}_z &= -i\hbar \frac{\partial}{\partial \phi} \label{eq:lz_spherical} \\
\hat{\mathbf{L}}^2 &= -\hbar^2 \left[ \frac{1}{\sin\theta} \frac{\partial}{\partial \theta}\left( \sin\theta \frac{\partial}{\partial \theta} \right) + \frac{1}{\sin^2\theta} \frac{\partial^2}{\partial \phi^2} \right] = -\hbar^2 \nabla_{\theta, \phi}^2 \label{eq:l2_spherical}
\end{align}

เมื่อเปรียบเทียบสมการ~\eqref{eq:l2_spherical} กับสมการค่าเจาะจงเชิงมุม~\eqref{eq:angular_eigen_eq} จะประจักษ์ชัดทันทีว่า
\begin{equation}
\hat{\mathbf{L}}^2 Y(\theta, \phi) = \hbar^2 \ell(\ell + 1) Y(\theta, \phi)
\end{equation}
นั่นคือ ฟังก์ชัน $Y(\theta, \phi)$ ก็คือสถานะเจาะจงของขนาดกำลังสองของโมเมนตัมเชิงมุมนั่นเอง

เราสามารถแยกตัวแปรของฟังก์ชันเชิงมุมต่อไปได้อีกขั้นหนึ่งเป็น
\begin{equation}
Y(\theta, \phi) = \Theta(\theta) \Phi(\phi)
\end{equation}
สมการค่าเจาะจงของ $\hat{L}_z$ คือ
\begin{equation}
\hat{L}_z Y(\theta, \phi) = -i\hbar \frac{\partial}{\partial \phi}\big[ \Theta(\theta)\Phi(\phi) \big] = m\hbar \Theta(\theta)\Phi(\phi)
\end{equation}
ซึ่งให้สมการอนุพันธ์ธรรมดาสำหรับ $\Phi(\phi)$
\begin{equation}
\frac{d\Phi}{d\phi} = i m \Phi(\phi) \quad \implies \quad \Phi_m(\phi) = \frac{1}{\sqrt{2\pi}} e^{im\phi}
\label{eq:phi_solution}
\end{equation}
เนื่องจากเมื่อมุม $\phi$ หมุนวนไปครบหนึ่งรอบ $2\pi$ พิกัดจะกลับมาที่ตำแหน่งเดิมในปริภูมิ ฟังก์ชันคลื่นจะต้องมีค่าเดี่ยวเชิงฟิสิกส์ (Single-valued Boundary Condition)
\begin{equation}
\Phi(\phi + 2\pi) = \Phi(\phi) \quad \implies \quad e^{im(\phi + 2\pi)} = e^{im\phi} \quad \implies \quad e^{i 2\pi m} = 1
\end{equation}
เงื่อนไขนี้บังคับให้ $m$ ต้องเป็นจำนวนเต็มเท่านั้น
\begin{equation}
m \in \{ 0, \pm 1, \pm 2, \dots \}
\end{equation}
โดยที่ $m$ เรียกว่า \textbf{เลขควอนตัมแม่เหล็ก (Magnetic Quantum Number)}

สำหรับส่วนของมุมขั้ว $\Theta(\theta)$ สมการที่ได้คือ \textbf{สมการพหุนามเลอฌ็องดร์คู่ควบ (Associated Legendre Differential Equation)} กำหนดตัวแปรใหม่ $x = \cos\theta$
\begin{equation}
(1 - x^2) \frac{d^2 \Theta}{dx^2} - 2x \frac{d\Theta}{dx} + \left[ \ell(\ell + 1) - \frac{m^2}{1 - x^2} \right] \Theta = 0
\label{eq:legendre_assoc_eq}
\end{equation}

ผลเฉลยที่ต่อเนื่องและมีค่าจำกัดที่ขั้วทั้งสอง ($\theta = 0$ และ $\theta = \pi$ หรือ $x = \pm 1$) จะเกิดขึ้นได้ก็ต่อเมื่อ $\ell$ เป็นจำนวนเต็มไม่เป็นลบ และ $|m| \le \ell$
\begin{equation}
\ell \in \{ 0, 1, 2, 3, \dots \}, \qquad m \in \{ -\ell, -\ell+1, \dots, 0, \dots, \ell-1, \ell \}
\end{equation}
สำหรับแต่ละค่า $\ell$ จะมีค่า $m$ ที่เป็นไปได้ทั้งหมด $2\ell + 1$ ค่า

ผลเฉลยของสมการ~\eqref{eq:legendre_assoc_eq} คือฟังก์ชันเลอฌ็องดร์คู่ควบ $P_\ell^m(x)$ ซึ่งหาได้จากสูตรของโรดริเกส (Rodrigues' Formula) สำหรับพหุนามเลอฌ็องดร์ $P_\ell(x)$
\begin{equation}
P_\ell(x) = \frac{1}{2^\ell \ell!} \frac{d^\ell}{dx^\ell} (x^2 - 1)^\ell
\end{equation}
\begin{equation}
P_\ell^m(x) = (-1)^m (1 - x^2)^{m/2} \frac{d^m}{dx^m} P_\ell(x) \quad (m \ge 0)
\end{equation}

ฟังก์ชันคลื่นเชิงมุมที่นอร์มอลไลซ์แล้วบนผิวทรงกลมหนึ่งหน่วยเรียกว่า \textbf{ฮาร์มอนิกส์ทรงกลม (Spherical Harmonics, $Y_\ell^m(\theta, \phi)$)}
\begin{equation}
Y_\ell^m(\theta, \phi) = (-1)^m \sqrt{\frac{2\ell + 1}{4\pi} \frac{(\ell - m)!}{(\ell + m)!}} \, P_\ell^m(\cos\theta) e^{im\phi} \quad (m \ge 0)
\label{eq:spherical_harmonics_exact}
\end{equation}
และสำหรับค่า $m$ ที่เป็นลบ กำหนดโดยสมมาตรสังยุคเชิงซ้อน
\begin{equation}
Y_\ell^{-m}(\theta, \phi) = (-1)^m \big[ Y_\ell^m(\theta, \phi) \big]^*
\end{equation}

\begin{table}[H]
\centering
\caption{ฟังก์ชันฮาร์มอนิกส์ทรงกลม $Y_\ell^m(\theta, \phi)$ สำหรับระดับ $\ell = 0, 1, 2$}
\label{tab:spherical_harmonics}
\begin{tabular}{cccl}
\hline
$\ell$ & $m$ & สัญลักษณ์ออร์บิทัล & ฟังก์ชันฮาร์มอนิกส์ทรงกลม $Y_\ell^m(\theta, \phi)$ \\
\hline
0 & 0 & $s$ & $Y_0^0 = \frac{1}{\sqrt{4\pi}}$ \\[6pt]
1 & 0 & $p_z$ & $Y_1^0 = \sqrt{\frac{3}{4\pi}} \cos\theta$ \\[6pt]
1 & $\pm 1$ & $p_x, p_y$ & $Y_1^{\pm 1} = \mp \sqrt{\frac{3}{8\pi}} \sin\theta \, e^{\pm i\phi}$ \\[6pt]
2 & 0 & $d_{z^2}$ & $Y_2^0 = \sqrt{\frac{5}{16\pi}} (3\cos^2\theta - 1)$ \\[6pt]
2 & $\pm 1$ & $d_{xz}, d_{yz}$ & $Y_2^{\pm 1} = \mp \sqrt{\frac{15}{8\pi}} \sin\theta\cos\theta \, e^{\pm i\phi}$ \\[6pt]
2 & $\pm 2$ & $d_{xy}, d_{x^2-y^2}$ & $Y_2^{\pm 2} = \sqrt{\frac{15}{32\pi}} \sin^2\theta \, e^{\pm 2i\phi}$ \\
\hline
\end{tabular}
\end{table}

ฮาร์มอนิกส์ทรงกลมมีสมบัติการตั้งฉากปรกติ (Orthonormality) บนผิวทรงกลมสองมิติ
\begin{equation}
\int_0^{2\pi} d\phi \int_0^\pi d\theta \sin\theta \, \big[ Y_\ell^m(\theta, \phi) \big]^* Y_{\ell'}^{m'}(\theta, \phi) = \delta_{\ell \ell'} \delta_{m m'}
\label{eq:spherical_harmonics_orthonormality}
\end{equation}
และภายใต้การสะท้อนพิกัดผ่านจุดกำเนิด (Parity Operation, $\mathbf{r} \to -\mathbf{r}$ ซึ่งเทียบเท่ากับ $\theta \to \pi - \theta$ และ $\phi \to \phi + \pi$) ฮาร์มอนิกส์ทรงกลมมีพาริตีที่แน่นอนเท่ากับ $(-1)^\ell$
\begin{equation}
\hat{\Pi} Y_\ell^m(\theta, \phi) = Y_\ell^m(\pi - \theta, \phi + \pi) = (-1)^\ell Y_\ell^m(\theta, \phi)
\end{equation}

% ==========================================================
% 6.3 อะตอมไฮโดรเจนและการแยกตัวแปรในแนวรัศมี
% ==========================================================
\section{อะตอมไฮโดรเจนและการแยกตัวแปรในแนวรัศมี}
\index{อะตอมไฮโดรเจน (Hydrogen Atom)}
\index{มวลลดทอน (Reduced Mass)}
\index{ศักย์ยังผล (Effective Potential)}

อะตอมไฮโดรเจนประกอบด้วยโปรตอนประจุบวก $+e$ มวล $M_p$ เป็นนิวเคลียส และอิเล็กตรอนประจุลบ $-e$ มวล $m_e$ โคจรอยู่โดยรอบ ระบบสองอนุภาคนี้สามารถแปลงเป็นระบบอนุภาคเดี่ยวได้อย่างแม่นตรงโดยการแยกการเคลื่อนที่ของจุดศูนย์กลางมวลออก และกำหนดให้มี \textbf{มวลลดทอน (Reduced Mass, $\mu$)}
\begin{equation}
\mu \equiv \frac{m_e M_p}{m_e + M_p} \approx m_e \left( 1 - \frac{m_e}{M_p} \right) \approx 0.999456 \, m_e
\label{eq:reduced_mass}
\end{equation}
เนื่องจากมวลโปรตอนมีค่ามากกว่าอิเล็กตรอนประมาณ 1836 เท่า มวลลดทอนจึงมีค่าใกล้เคียงกับมวลของอิเล็กตรอนมาก แต่การใช้มวลลดทอนช่วยให้อธิบายสเปกตรัมของไอโซโทป เช่น ดิวเทอเรียม ($^2\text{H}$) และโพซิตโรเนียมได้อย่างแม่นยำ

สำหรับไอออนที่มีอิเล็กตรอนตัวเดียวและมีโปรตอน $Z$ ตัวในสนามศักย์คูลอมบ์ (Hydrogen-like ions เช่น $\text{He}^+, \text{Li}^{2+}$) พลังงานศักย์คือ
\begin{equation}
V(r) = -\frac{Z e^2}{4\pi\varepsilon_0 r}
\label{eq:coulomb_potential_z}
\end{equation}

นำศักย์คูลอมบ์แทนลงในสมการเชิงรัศมี~\eqref{eq:radial_general_eq}
\begin{equation}
\frac{1}{r^2}\frac{d}{dr}\left( r^2 \frac{dR}{dr} \right) + \frac{2\mu}{\hbar^2}\left[ E + \frac{Ze^2}{4\pi\varepsilon_0 r} - \frac{\hbar^2 \ell(\ell + 1)}{2\mu r^2} \right] R(r) = 0
\label{eq:radial_differential_full}
\end{equation}

เพื่อแปลงสมการให้อยู่ในรูปที่คุ้นเคย เรานิยามฟังก์ชันคลื่นรัศมีปรับปรุง $u(r) \equiv r R(r)$ เมื่อคำนวณอนุพันธ์
\begin{equation}
\frac{d}{dr}\left( r^2 \frac{dR}{dr} \right) = r \frac{d^2 u}{dr^2}
\end{equation}
แทนลงในสมการ~\eqref{eq:radial_differential_full} จะได้สมการที่มีรูปแบบเทียบเท่าสมการชเรอดิงเงอร์หนึ่งมิติ
\begin{equation}
-\frac{\hbar^2}{2\mu} \frac{d^2 u}{dr^2} + V_{\text{eff}}(r) u(r) = E u(r)
\label{eq:radial_schrodinger_u}
\end{equation}
โดยที่ \textbf{ศักย์ยังผล (Effective Potential, $V_{\text{eff}}$)} ประกอบด้วยสองส่วน
\begin{equation}
V_{\text{eff}}(r) = -\frac{Ze^2}{4\pi\varepsilon_0 r} + \frac{\hbar^2 \ell(\ell + 1)}{2\mu r^2}
\label{eq:effective_potential_def}
\end{equation}
ส่วนแรกคือแรงดึงดูดคูลอมบ์ระหว่างประจุต่างชนิด และส่วนหลังคือ \textbf{กำแพงหนีศูนย์กลางควอนตัม (Centrifugal Barrier)} ซึ่งทำหน้าที่เสมือนแรงผลักที่ป้องกันไม่ให้อนุภาคที่มีโมเมนตัมเชิงมุมเข้าใกล้จุดศูนย์กลาง สำหรับสถานะที่มี $\ell > 0$ เมื่อ $r \to 0$ เทอมหนีศูนย์กลางจะครอบงำ ส่งผลให้ความน่าจะเป็นในการพบอนุภาคที่จุดศูนย์กลางลดลงเป็นศูนย์

\subsection{พฤติกรรมเชิงเส้นกำกับและการแปลงตัวแปรไร้มิติ}
สำหรับสถานะผูกพัน (Bound States) พลังงานรวมมีค่าเป็นลบ $E < 0$ เรานิยามพารามิเตอร์คงตัว
\begin{equation}
\kappa \equiv \frac{\sqrt{-2\mu E}}{\hbar} > 0
\end{equation}
และตัวแปรระยะทางไร้มิติ
\begin{equation}
\rho \equiv 2\kappa r
\end{equation}
และนิยามค่าคงตัวไร้มิติ
\begin{equation}
\rho_0 \equiv \frac{2\mu Z e^2}{4\pi\varepsilon_0 \hbar^2 \kappa} = \frac{2 Z}{\kappa a_0^*}
\end{equation}
โดยที่ $a_0^* = \frac{4\pi\varepsilon_0 \hbar^2}{\mu e^2}$ คือรัศมีโบร์ปรับปรุงด้วยมวลลดทอน

เมื่อแทนค่าตัวแปรไร้มิติลงในสมการ~\eqref{eq:radial_schrodinger_u} จะได้
\begin{equation}
\frac{d^2 u}{d\rho^2} = \left[ \frac{1}{4} - \frac{\rho_0}{\rho} + \frac{\ell(\ell + 1)}{\rho^2} \right] u(\rho)
\label{eq:dimensionless_radial_eq}
\end{equation}

เราวิเคราะห์พฤติกรรมของสมการที่ขอบเขตปลายทั้งสองด้าน
\begin{enumerate}
    \item \textbf{ที่ระยะอนันต์ ($\rho \to \infty$)} ---- เทอม $-\rho_0/\rho$ และ $\ell(\ell+1)/\rho^2$ เข้าใกล้ศูนย์ สมการลดรูปเป็น $\frac{d^2 u}{d\rho^2} \approx \frac{1}{4} u$ ซึ่งมีผลเฉลยทั่วไป $u(\rho) \sim A e^{-\rho/2} + B e^{\rho/2}$ เพื่อให้ฟังก์ชันคลื่นลู่เข้าและสามารถนอร์มอลไลซ์ได้ เราต้องกำหนด $B = 0$ นั่นคือ $u(\rho) \sim e^{-\rho/2}$
    \item \textbf{ที่จุดกำเนิด ($\rho \to 0$)} ---- เทอม $\ell(\ell+1)/\rho^2$ มีค่าครอบงำที่สุด สมการลดรูปเป็น $\frac{d^2 u}{d\rho^2} \approx \frac{\ell(\ell+1)}{\rho^2} u$ สมมติผลเฉลย $u \sim \rho^s$ จะได้ $s(s-1) = \ell(\ell+1)$ ซึ่งให้ค่า $s = \ell + 1$ หรือ $s = -\ell$ ค่า $s = -\ell$ ทำให้ฟังก์ชันคลื่นระเบิดลู่ออกที่จุดกำเนิดและทำให้พลังงานจลน์มีค่าเป็นอนันต์ จึงต้องตัดทิ้ง เหลือเพียง $u(\rho) \sim \rho^{\ell + 1}$
\end{enumerate}

เมื่อสกัดพฤติกรรมที่จุดขอบเขตทั้งสองด้านออกแล้ว เราสามารถเขียนฟังก์ชันคลื่นในรูปผลคูณ
\begin{equation}
u(\rho) = \rho^{\ell + 1} e^{-\rho/2} v(\rho)
\label{eq:ansatz_u}
\end{equation}
เมื่อนำสมการ~\eqref{eq:ansatz_u} แทนกลับลงในสมการ~\eqref{eq:dimensionless_radial_eq} จะได้สมการอนุพันธ์สำหรับฟังก์ชัน $v(\rho)$
\begin{equation}
\rho \frac{d^2 v}{d\rho^2} + \big[ 2(\ell + 1) - \rho \big] \frac{dv}{d\rho} + \left( \rho_0 - \ell - 1 \right) v = 0
\label{eq:kummers_equation}
\end{equation}
สมการ~\eqref{eq:kummers_equation} คือ \textbf{สมการไฮเพอร์จีโอเมตริกคอนฟลูเอนต์ (Confluent Hypergeometric Equation)} หรือสมการลาแกร์คู่ควบ

\subsection{การกระจายอนุกรมกำลังและการควอนไตซ์พลังงาน}
เราแก้สมการ~\eqref{eq:kummers_equation} ด้วยระเบียบวิธีอนุกรมกำลังของโฟรเบนิอุส (Frobenius Method)
\begin{equation}
v(\rho) = \sum_{j=0}^\infty a_j \rho^j
\end{equation}
หาอนุพันธ์อันดับหนึ่งและอันดับสอง
\begin{equation}
\frac{dv}{d\rho} = \sum_{j=0}^\infty j a_j \rho^{j-1}, \qquad \frac{d^2 v}{d\rho^2} = \sum_{j=0}^\infty j(j-1) a_j \rho^{j-2}
\end{equation}
แทนลงในสมการ~\eqref{eq:kummers_equation} และจัดสัมประสิทธิ์ของกำลัง $\rho^j$ ให้เท่ากับศูนย์ จะได้ \textbf{ความสัมพันธ์เวียนเกิด (Recurrence Relation)}
\begin{equation}
a_{j+1} = \frac{j + \ell + 1 - \rho_0}{(j + 1)(j + 2\ell + 2)} a_j
\label{eq:recurrence_relation}
\end{equation}

เมื่อพิจารณาที่ $j$ มีค่ามากๆ ($j \to \infty$) อัตราส่วนระหว่างพจน์ที่อยู่ติดกันคือ
\begin{equation}
\frac{a_{j+1}}{a_j} \approx \frac{j}{j^2} = \frac{1}{j}
\end{equation}
ซึ่งเป็นพฤติกรรมเดียวกันทุกประการกับการกระจายอนุกรมของฟังก์ชันเอกซ์โพเนนเชียล $e^\rho = \sum \rho^j/j!$ ดังนั้น หากอนุกรมนี้ไม่สิ้นสุดลงที่พจน์จำกัด ฟังก์ชัน $v(\rho)$ จะประพฤติตนเป็น $e^\rho$ ที่ระยะไกล ซึ่งจะทำให้
\begin{equation}
u(\rho) \sim \rho^{\ell+1} e^{-\rho/2} e^\rho = \rho^{\ell+1} e^{+\rho/2} \xrightarrow{\rho \to \infty} \infty
\end{equation}
ฟังก์ชันคลื่นจะลู่ออกสู่อนันต์และไม่สามารถนอร์มอลไลซ์ได้ในเชิงฟิสิกส์

ดังนั้น \textbf{อนุกรมกำลังจะต้องสิ้นสุดลงที่พจน์จำกัด (Polynomial Termination)} สมมติว่ามีดัชนีสูงสุด $j_{\text{max}} = N$ ที่ทำให้ $a_{N+1} = 0$ จากสมการ~\eqref{eq:recurrence_relation} ตัวเศษจะต้องมีค่าเป็นศูนย์
\begin{equation}
N + \ell + 1 - \rho_0 = 0 \quad \implies \quad \rho_0 = N + \ell + 1
\end{equation}
เนื่องจาก $N \in \{0, 1, 2, \dots\}$ และ $\ell \in \{0, 1, 2, \dots\}$ ผลรวม $N + \ell + 1$ จึงเป็นจำนวนเต็มบวกเสมอ เรานิยามจำนวนเต็มนี้เป็น \textbf{เลขควอนตัมหลัก (Principal Quantum Number, $n$)}
\begin{equation}
n \equiv N + \ell + 1 \in \{ 1, 2, 3, 4, \dots \}
\label{eq:principal_quantum_number}
\end{equation}
จากสมการ~\eqref{eq:principal_quantum_number} เนื่องจาก $N \ge 0$ สำหรับค่า $n$ ที่กำหนด ค่าของเลขควอนตัมออร์บิทัล $\ell$ จึงถูกจำกัดโดยอัตโนมัติ
\begin{equation}
\ell \in \{ 0, 1, 2, \dots, n - 1 \}
\end{equation}

จากความสัมพันธ์ $\rho_0 = n$ เราแทนค่ากลับไปยังพลังงาน $E$
\begin{equation}
\rho_0 = \frac{2\mu Z e^2}{4\pi\varepsilon_0 \hbar^2 \kappa} = n \quad \implies \quad \kappa = \frac{\mu Z e^2}{4\pi\varepsilon_0 \hbar^2 n}
\end{equation}
เมื่อแทน $\kappa^2 = -2\mu E/\hbar^2$ เราได้สูตร \textbf{ระดับพลังงานควอนไตซ์ของอะตอมไฮโดรเจนและไฮโดรเจนิกไอออน}
\begin{equation}
E_n = -\frac{\mu Z^2 e^4}{32\pi^2 \varepsilon_0^2 \hbar^2} \frac{1}{n^2} = -\frac{Z^2 E_R}{n^2} = -\frac{13.606 \, Z^2}{n^2}~\text{eV}
\label{eq:hydrogen_energy_formula}
\end{equation}
โดยที่ $E_R \equiv \frac{\mu e^4}{32\pi^2\varepsilon_0^2\hbar^2} \approx 13.606~\text{eV}$ คือพลังงานริดเบิร์ก (Rydberg Energy) และมาตรวัดระยะทางธรรมชาติของอะตอมคือกำหนดโดย \textbf{รัศมีโบร์ (Bohr Radius)}
\begin{equation}
a_0 \equiv \frac{4\pi\varepsilon_0 \hbar^2}{m_e e^2} \approx 0.529177 \times 10^{-10}~\text{m} = 0.529~\text{\AA}
\end{equation}

\subsection{พหุนามลาแกร์คู่ควบและฟังก์ชันคลื่นเชิงรัศมี}
ผลเฉลยพหุนามจำกัดของสมการ~\eqref{eq:kummers_equation} สามารถเขียนในรูปของ \textbf{พหุนามลาแกร์คู่ควบ (Associated Laguerre Polynomials, $L_p^k(x)$)}
\begin{equation}
v(\rho) = L_{n - \ell - 1}^{2\ell + 1}(\rho)
\end{equation}
โดยที่พหุนามลาแกร์คู่ควบนิยามผ่านสูตรของโรดริเกส
\begin{equation}
L_p^k(x) = (-1)^k \frac{d^k}{dx^k} \big[ L_{p+k}(x) \big] = \frac{x^{-k} e^x}{p!} \frac{d^p}{dx^p}\left( x^{p+k} e^{-x} \right)
\end{equation}

เมื่อรวมทุกเทอมและทำการนอร์มอลไลซ์ $\int_0^\infty |R_{n\ell}(r)|^2 r^2 dr = 1$ ฟังก์ชันคลื่นเชิงรัศมีที่สมบูรณ์มีรูปดังนี้
\begin{equation}
R_{n\ell}(r) = -\sqrt{\left(\frac{2Z}{n a_0}\right)^3 \frac{(n - \ell - 1)!}{2n \big[(n + \ell)!\big]^3}} \, e^{-\rho/2} \, \rho^\ell \, L_{n - \ell - 1}^{2\ell + 1}(\rho)
\label{eq:radial_wavefunction_general}
\end{equation}
โดยที่ $\rho = \frac{2Z r}{n a_0}$

\begin{table}[H]
\centering
\caption{ฟังก์ชันคลื่นเชิงรัศมี $R_{n\ell}(r)$ ของอะตอมไฮโดรเจน ($Z = 1$)}
\label{tab:radial_wavefunctions}
\begin{tabular}{ccl}
\hline
$n$ & $\ell$ & ฟังก์ชันคลื่นเชิงรัศมี $R_{n\ell}(r)$ \\
\hline
1 & 0 ($1s$) & $R_{10}(r) = 2 a_0^{-3/2} e^{-r/a_0}$ \\[6pt]
2 & 0 ($2s$) & $R_{20}(r) = \frac{1}{2\sqrt{2}} a_0^{-3/2} \left( 2 - \frac{r}{a_0} \right) e^{-r/2a_0}$ \\[6pt]
2 & 1 ($2p$) & $R_{21}(r) = \frac{1}{2\sqrt{6}} a_0^{-3/2} \left( \frac{r}{a_0} \right) e^{-r/2a_0}$ \\[6pt]
3 & 0 ($3s$) & $R_{30}(r) = \frac{2}{81\sqrt{3}} a_0^{-3/2} \left( 27 - 18\frac{r}{a_0} + 2\frac{r^2}{a_0^2} \right) e^{-r/3a_0}$ \\[6pt]
3 & 1 ($3p$) & $R_{31}(r) = \frac{4}{81\sqrt{6}} a_0^{-3/2} \left( 6\frac{r}{a_0} - \frac{r^2}{a_0^2} \right) e^{-r/3a_0}$ \\[6pt]
3 & 2 ($3d$) & $R_{32}(r) = \frac{4}{81\sqrt{30}} a_0^{-3/2} \left( \frac{r}{a_0} \right)^2 e^{-r/3a_0}$ \\
\hline
\end{tabular}
\end{table}

\subsection{ภาวะเสื่อมของระดับพลังงานและเวกเตอร์รุนเก-เลนซ์}
พิจารณาสูตรพลังงาน~\eqref{eq:hydrogen_energy_formula} จะสังเกตเห็นว่า พลังงาน $E_n$ ขึ้นอยู่กับเลขควอนตัมหลัก $n$ เพียงตัวเดียว โดยไม่ขึ้นกับเลขควอนตัมออร์บิทัล $\ell$ หรือเลขควอนตัมแม่เหล็ก $m$ เลย สถานการณ์นี้เรียกว่า \textbf{ภาวะเสื่อม (Degeneracy)}

สำหรับแต่ละค่า $n$ ที่กำหนด
\begin{itemize}
    \item ค่า $\ell$ มีได้ตั้งแต่ $0$ ถึง $n - 1$
    \item สำหรับแต่ละค่า $\ell$ มีค่า $m$ ที่เป็นไปได้ทั้งหมด $2\ell + 1$ ค่า
\end{itemize}
ดังนั้น ระดับความเสื่อมเชิงพื้นที่รวม (Total Spatial Degeneracy, $g_n$) คือ
\begin{equation}
g_n = \sum_{\ell = 0}^{n - 1} (2\ell + 1) = 2 \left( \frac{(n - 1)n}{2} \right) + n = n(n - 1) + n = n^2
\label{eq:spatial_degeneracy}
\end{equation}
เมื่อนำสปินของอิเล็กตรอนมาร่วมพิจารณา ($m_s = \pm 1/2$) ความเสื่อมรวมทั้งหมดจะมีค่าเป็น $2n^2$

ความเสื่อมตามเลขควอนตัม $m$ เป็นผลมาจากสมมาตรการหมุนของระบบสามมิติทั่วไป (Rotational Symmetry $SO(3)$) ทว่าความเสื่อมตามเลขควอนตัม $\ell$ ที่พลังงานของสถานะ $2s$ และ $2p$ เท่ากันพอดีในศักย์ $1/r$ เรียกว่า \textbf{ความเสื่อมโดยบังเอิญ (Accidental Degeneracy)} ซึ่งแท้จริงแล้วไม่ใช่เรื่องบังเอิญ แต่เกิดจากการมีอยู่ของปริมาณคงตัวเชิงพลวัตอีกตัวหนึ่ง คือ \textbf{เวกเตอร์ลาปลาซ-รุนเก-เลนซ์ (Laplace-Runge-Lenz Vector, $\hat{\mathbf{A}}$)}
\begin{equation}
\hat{\mathbf{A}} = \frac{1}{2\mu} (\hat{\mathbf{p}} \times \hat{\mathbf{L}} - \hat{\mathbf{L}} \times \hat{\mathbf{p}}) - \frac{Ze^2}{4\pi\varepsilon_0} \frac{\mathbf{r}}{r}
\end{equation}
ซึ่งสลับที่กับฮามิลโทเนียน $[\hat{\mathbf{A}}, \hat{H}] = 0$ ส่งผลให้ระบบอะตอมไฮโดรเจนมีสมมาตรซ่อนเร้นระดับสูงในสี่มิติ คือกลุ่มสมมาตร $SO(4)$ ซึ่งเป็นสาเหตุทางคณิตศาสตร์ที่แท้จริงของความเสื่อม $n^2$

% ==========================================================
% 6.4 รูปร่างของออร์บิทัล ความน่าจะเป็น และค่าคาดหวัง
% ==========================================================
\section{รูปร่างของออร์บิทัล ความน่าจะเป็น และค่าคาดหวัง}
\index{ออร์บิทัลเชิงอะตอม (Atomic Orbitals)}
\index{ความหนาแน่นความน่าจะเป็นเชิงรัศมี (Radial Probability Density)}
\index{บัพเชิงรัศมี (Radial Nodes)}

ฟังก์ชันคลื่นรวมของสถานะคงตัวในอะตอมไฮโดรเจนเขียนได้ในรูป
\begin{equation}
\psi_{n\ell m}(r, \theta, \phi) = R_{n\ell}(r) Y_\ell^m(\theta, \phi)
\end{equation}
ความน่าจะเป็นในการค้นพบอิเล็กตรอนในองค์ประกอบปริมาตร $d^3\mathbf{r} = r^2 \sin\theta \, dr d\theta d\phi$ คือ
\begin{equation}
dP = |\psi_{n\ell m}(r, \theta, \phi)|^2 d^3\mathbf{r} = |R_{n\ell}(r)|^2 |Y_\ell^m(\theta, \phi)|^2 r^2 \sin\theta \, dr d\theta d\phi
\end{equation}

หากเราสนใจเฉพาะโอกาสที่จะพบอิเล็กตรอนที่ระยะห่างระหว่าง $r$ ถึง $r + dr$ โดยไม่คำนึงถึงทิศทางมุมรอบจุดกำเนิด เราทำการอินทิเกรตฟังก์ชันคลื่นบนมุมเชิงทรงกลมทั้งหมด
\begin{equation}
P_{n\ell}(r) dr = \left( \int_0^{2\pi} d\phi \int_0^\pi \sin\theta d\theta \, |Y_\ell^m(\theta, \phi)|^2 \right) |R_{n\ell}(r)|^2 r^2 dr
\end{equation}
เนื่องจากฮาร์มอนิกส์ทรงกลมนอร์มอลไลซ์บนผิวทรงกลมแล้ว อินทิกรัลมุมจึงมีค่าเท่ากับ 1 พอดี ทำให้ได้ \textbf{ความหนาแน่นความน่าจะเป็นในแนวรัศมี (Radial Probability Density)}
\begin{equation}
P_{n\ell}(r) = r^2 |R_{n\ell}(r)|^2
\label{eq:radial_probability_density_formula}
\end{equation}

\subsection{บัพเชิงรัศมีและบัพเชิงมุม}
จำนวนตำแหน่งที่ฟังก์ชันคลื่นมีค่าเป็นศูนย์ เรียกว่า \textbf{บัพ (Nodes)} ซึ่งแบ่งออกเป็นสองประเภท
\begin{enumerate}
    \item \textbf{บัพเชิงรัศมี (Radial Nodes)} ---- จุดที่ $R_{n\ell}(r) = 0$ (ไม่นับที่ $r = 0$ และ $r = \infty$) กำหนดโดยระดับขั้นของพหุนามลาแกร์คู่ควบ
    \begin{equation}
    N_{\text{radial}} = n - \ell - 1
    \end{equation}
    \item \textbf{บัพเชิงมุม (Angular Nodes)} ---- พื้นผิวทรงกรวยหรือระนาบที่ $Y_\ell^m(\theta, \phi) = 0$ กำหนดโดยเลขควอนตัมออร์บิทัล
    \begin{equation}
    N_{\text{angular}} = \ell
    \end{equation}
\end{enumerate}
ดังนั้น จำนวนบัพรวมทั้งหมดของออร์บิทัลคือ
\begin{equation}
N_{\text{total}} = N_{\text{radial}} + N_{\text{angular}} = (n - \ell - 1) + \ell = n - 1
\end{equation}

ตัวอย่างเช่น
\begin{itemize}
    \item สถานะ $1s$ ($n=1, \ell=0$) ---- บัพเชิงรัศมี 0, บัพเชิงมุม 0, รวม 0 บัพ
    \item สถานะ $2s$ ($n=2, \ell=0$) ---- บัพเชิงรัศมี 1, บัพเชิงมุม 0, รวม 1 บัพ ที่ระยะ $r = 2a_0$
    \item สถานะ $2p$ ($n=2, \ell=1$) ---- บัพเชิงรัศมี 0, บัพเชิงมุม 1 (ระนาบกั้นระหว่างพู), รวม 1 บัพ
    \item สถานะ $3p$ ($n=3, \ell=1$) ---- บัพเชิงรัศมี 1, บัพเชิงมุม 1, รวม 2 บัพ
\end{itemize}

\subsection{ค่าคาดหวังของระยะทางและทฤษฎีบทไวเรียล}
การคำนวณค่าคาดหวังของกำลังต่างๆ ของระยะทาง $\langle r^k \rangle$ สามารถหาได้จากการอินทิเกรตฟังก์ชันคลื่นรัศมี หรือใช้ความสัมพันธ์เวียนเกิดของคราเมอส์ (Kramers' Recurrence Relation) สำหรับอะตอมไฮโดรเจน ($Z = 1$)
\begin{align}
\langle r \rangle_{n\ell} &= \frac{a_0}{2} \Big[ 3n^2 - \ell(\ell + 1) \Big] \label{eq:exp_r} \\
\langle r^2 \rangle_{n\ell} &= \frac{a_0^2 n^2}{2} \Big[ 5n^2 + 1 - 3\ell(\ell + 1) \Big] \label{eq:exp_r2} \\
\left\langle \frac{1}{r} \right\rangle_{n\ell} &= \frac{1}{n^2 a_0} \label{eq:exp_inv_r} \\
\left\langle \frac{1}{r^2} \right\rangle_{n\ell} &= \frac{1}{n^3 (\ell + 1/2) a_0^2} \label{eq:exp_inv_r2} \\
\left\langle \frac{1}{r^3} \right\rangle_{n\ell} &= \frac{1}{n^3 \ell(\ell + 1/2)(\ell + 1) a_0^3} \quad (\ell > 0) \label{eq:exp_inv_r3}
\end{align}

จากสมการ~\eqref{eq:exp_inv_r} เราคำนวณค่าคาดหวังของพลังงานศักย์ไฟฟ้าคูลอมบ์
\begin{equation}
\langle V \rangle = -\frac{e^2}{4\pi\varepsilon_0} \left\langle \frac{1}{r} \right\rangle = -\frac{e^2}{4\pi\varepsilon_0 n^2 a_0} = 2 E_n
\end{equation}
และเนื่องจาก $E_n = \langle K \rangle + \langle V \rangle$ จะได้
\begin{equation}
\langle K \rangle = -E_n = -\frac{1}{2}\langle V \rangle
\end{equation}
ความสัมพันธ์ $2\langle K \rangle + \langle V \rangle = 0$ นี้คือ \textbf{ทฤษฎีบทไวเรียล (Virial Theorem)} สำหรับศักย์ผกผันกำลังหนึ่ง ซึ่งพิสูจน์ความสอดคล้องอย่างสมบูรณ์แบบระหว่างกลศาสตร์ควอนตัมและกลศาสตร์ดั้งเดิม

% ==========================================================
% 6.5 การทดลองชแตร์น-แกร์ลัคและสปินของอิเล็กตรอน
% ==========================================================
\section{การทดลองชแตร์น-แกร์ลัคและสปินของอิเล็กตรอน}
\index{การทดลองชแตร์น-แกร์ลัค (Stern-Gerlach Experiment)}
\index{สปินของอิเล็กตรอน (Electron Spin)}
\index{โบร์แมกเนตอน (Bohr Magneton)}

ในคริสต์ศักราช 1922 ณ สถาบันฟิสิกส์แห่งแฟรงก์เฟิร์ต ออตโต ชแตร์น (Otto Stern, 1888--1969) และวัลเทอร์ แกร์ลัค (Walther Gerlach, 1889--1979) ได้ออกแบบการทดลองครั้งประวัติศาสตร์เพื่อทดสอบสมมติฐานการควอนไตซ์เชิงทิศทางในปริภูมิ (Space Quantization) ของแบบจำลองอะตอมของบอร์-ซอมเมอร์เฟลด์

\subsection{กลไกเชิงฟิสิกส์ของการทดลอง}
ในการทดลอง ลำอะตอมของธาตุเงิน (Silver, Ag) บริสุทธิ์ถูกทำให้ระเหยในเตาอบสุญญากาศอุณหภูมิสูง ลำอะตอมที่เป็นกลางทางไฟฟ้าวิ่งผ่านช่องคอลลิเมเตอร์ให้เป็นลำขนานแคบ แล้วพุ่งผ่านบริเวณที่มี \textbf{สนามแม่เหล็กที่มีความไม่สม่ำเสมอสูง (Inhomogeneous Magnetic Field)} ในแนวแกน $z$

อะตอมเงินมีเลขอะตอม $Z = 47$ โดยมีการจัดเรียงอิเล็กตรอนเป็น $[\text{Kr}] 4d^{10} 5s^1$ อิเล็กตรอน 46 ตัวในชั้นในจับคู่กันเต็มทุกวงโคจร ส่วนอิเล็กตรอนตัวนอกสุดตัวที่ 47 อยู่ในสถานะ $5s$ ซึ่งมีโมเมนตัมเชิงมุมออร์บิทัลเป็นศูนย์ ($\ell = 0$) ดังนั้น ตามทฤษฎีฟิสิกส์ดั้งเดิมและกลศาสตร์ควอนตัมยุคแรก โมเมนต์แม่เหล็กออร์บิทัลของอะตอมเงินย่อมเป็นศูนย์อย่างสิ้นเชิง
\begin{equation}
\boldsymbol{\mu}_L = -\frac{e}{2m_e}\mathbf{L} = \mathbf{0}
\end{equation}

หากอะตอมมีโมเมนต์แม่เหล็กสุทธิ $\boldsymbol{\mu}$ พลังงานศักย์ของขั้วคู่แม่เหล็กในสนามแม่เหล็กภายนอก $\mathbf{B}$ คือ $U = -\boldsymbol{\mu}\cdot\mathbf{B}$ แรงสุทธิที่กระทำต่ออะตอมหาได้จากเกรเดียนต์ของพลังงานศักย์
\begin{equation}
\mathbf{F} = -\boldsymbol{\nabla} U = \boldsymbol{\nabla}(\boldsymbol{\mu}\cdot\mathbf{B})
\end{equation}
ขั้วแม่เหล็กได้รับการออกแบบให้ขั้วเหนือเป็นสันคมแหลมและขั้วใต้เป็นร่องเว้า ทำให้สนามแม่เหล็กมีความเข้มและเกรเดียนต์หลักตามแนวแกน $z$ ($\partial B_z/\partial z \gg \partial B_x/\partial x, \partial B_y/\partial y$)
\begin{equation}
F_z \approx \mu_z \frac{\partial B_z}{\partial z}
\label{eq:stern_gerlach_force}
\end{equation}

\begin{figure}[H]
\centering
\begin{tikzpicture}[scale=1.05, >=Stealth]
    % Oven source
    \draw[fill=slateGrey!30, draw=darkNavy, thick] (-4.5, -0.9) rectangle (-3.2, 0.9);
    \node at (-3.85, 0) [font=\tiny\bfseries, align=center] {เตาอบ\\อะตอมเงิน\\$T \approx 1300\text{ K}$};
    \draw[->, thick, darkNavy] (-3.2, 0) -- (-2.2, 0);

    % Collimator slits
    \draw[line width=3pt, slateGrey] (-2.2, -1.3) -- (-2.2, -0.2);
    \draw[line width=3pt, slateGrey] (-2.2, 0.2) -- (-2.2, 1.3);

    % Inhomogeneous Magnet Pole Pieces
    % Top sharp pole (N)
    \draw[fill=cyberBlue!25, draw=cyberBlue, line width=1.2pt] (-1.1, 1.8) -- (1.1, 1.8) -- (0, 0.6) -- cycle;
    \node at (0, 1.3) [font=\small\bfseries\color{darkNavy}] {ขั้ว N (ปลายแหลมคม)};

    % Bottom flat pole (S)
    \draw[fill=entanglePink!25, draw=entanglePink, line width=1.2pt] (-1.3, -1.8) rectangle (1.3, -0.9);
    \node at (0, -1.35) [font=\small\bfseries\color{darkNavy}] {ขั้ว S (ผิวร่องเรียบ)};

    % B-field gradient arrows
    \draw[->, dashed, slateGrey, line width=0.9pt] (0, 0.5) -- (0, -0.8) node[midway, right=3pt, font=\scriptsize] {$\frac{\partial B_z}{\partial z} > 0$};

    % Beam splitting
    \draw[line width=1.5pt, quantumGold] (-2.2, 0) -- (0, 0);
    \draw[->, line width=1.5pt, cyberBlue] (0, 0) -- (3.6, 1.0) node[right, font=\small\bfseries] {$m_s = +1/2$ (สปินชี้ขึ้น)};
    \draw[->, line width=1.5pt, entanglePink] (0, 0) -- (3.6, -1.0) node[right, font=\small\bfseries] {$m_s = -1/2$ (สปินชี้ลง)};

    % Classical prediction smudge (dashed)
    \draw[dashed, slateGrey!70, line width=0.8pt] (0,0) -- (3.6, 0);

    % Detector screen
    \draw[fill=slateGrey!15, draw=slateGrey, thick] (3.6, -1.8) rectangle (3.8, 1.8);
    \fill[cyberBlue] (3.7, 1.0) ellipse (2.0pt and 7pt);
    \fill[entanglePink] (3.7, -1.0) ellipse (2.0pt and 7pt);
    \node at (4.5, 0) [font=\tiny, align=left] {ผลการทดลอง\\แยกเป็น 2 แถบ\\ชัดเจน};
\end{tikzpicture}
\caption{แผนภาพการทดลองชแตร์น-แกร์ลัค แสดงการแยกตัวของลำอะตอมเงินในสนามแม่เหล็กไม่สม่ำเสมอเป็นสองแถบที่สมมาตร}
\label{fig:stern_gerlach_detailed}
\end{figure}

\subsection{การทำนายคลาสสิกเทียบกับผลการทดลองควอนตัม}
\begin{itemize}
    \item \textbf{ตามทฤษฎีคลาสสิก} ---- เวกเตอร์โมเมนต์แม่เหล็ก $\boldsymbol{\mu}$ ของอะตอมควรวางตัวแบบสุ่มในทุกทิศทางในทรงกลมสามมิติ องค์ประกอบ $\mu_z = |\boldsymbol{\mu}|\cos\theta$ จึงควรมีค่าต่อเนื่องได้ตั้งแต่ $-|\boldsymbol{\mu}|$ ถึง $+|\boldsymbol{\mu}|$ ลำอะตอมที่ตกกระทบฉากควรแผ่กระจายเป็นแถบปื้นต่อเนื่อง (Continuous Smudge) รอบจุดกึ่งกลาง
    \item \textbf{ผลการทดลองจริง} ---- ลำอะตอมเงินกลับแยกออกเป็น \textbf{สองแถบที่แยกจากกันอย่างเด็ดขาดและสมมาตร} โดยไม่มีอะตอมตกลงที่บริเวณจุดกึ่งกลางเลยแม้แต่ตัวเดียว!
\end{itemize}

การค้นพบนี้สร้างความฉงนแก่วงการฟิสิกส์อยู่หลายปี จนกระทั่งในคริสต์ศักราช 1925 สองนักฟิสิกส์หนุ่มชาวดัตช์ ซามูเอล เกาด์สมิต (Samuel Goudsmit, 1902--1978) และจอร์จ อูห์เลนเบก (George Uhlenbeck, 1900--1988) ได้เสนอแนวคิดปฏิวัติว่า อิเล็กตรอนมี \textbf{โมเมนตัมเชิงมุมสปิน (Spin Angular Momentum, $\mathbf{S}$)} ซึ่งเป็นสมบัติควอนตัมเฉพาะตัวภายใน (Intrinsic Property) เช่นเดียวกับมวลและประจุไฟฟ้า

โมเมนต์แม่เหล็กสปินของอิเล็กตรอนสัมพันธ์กับสปินผ่านอัตราส่วนไจโรแมกเนติก (Gyromagnetic Ratio)
\begin{equation}
\boldsymbol{\mu}_s = -g_s \frac{e}{2m_e} \mathbf{S} = -g_s \frac{\mu_B}{\hbar} \mathbf{S}
\label{eq:spin_magnetic_moment}
\end{equation}
โดยที่ $\mu_B \equiv \frac{e\hbar}{2m_e} \approx 9.274 \times 10^{-24}~\text{J/T}$ คือ \textbf{โบร์แมกเนตอน (Bohr Magneton)} และ $g_s \approx 2$ คือ $g$-factor ของสปินอิเล็กตรอน (ซึ่งต่อมาทฤษฎีควอนตัมอิเล็กโทรไดนามิกส์คำนวณได้ $g_s \approx 2.002319$)

เนื่องจากลำอะตอมแยกออกเป็น 2 แถบ จำนวนสถานะการฉายภาพตามแกน $z$ คือ $2s + 1 = 2$ ซึ่งสรุปได้โดยตรงว่า เลขควอนตัมสปินของอิเล็กตรอนคือ
\begin{equation}
s = \frac{1}{2}
\end{equation}
และค่าองค์ประกอบของสปินตามแนวแกน $z$ ควอนไตซ์เป็นสองค่าเท่านั้น
\begin{equation}
S_z = m_s \hbar = \pm \frac{1}{2}\hbar
\end{equation}

\subsection{การวัดชแตร์น-แกร์ลัคแบบเรียงลำดับ}
การทดลองชแตร์น-แกร์ลัคแบบต่อเนื่องหลายขั้นตอน (Sequential Stern-Gerlach Experiments) เป็นตัวอย่างเชิงประจักษ์ที่ทรงพลังที่สุดในการสาธิต \textbf{สัจพจน์การยุบตัวของฟังก์ชันคลื่น (Wavefunction Collapse)} และธรรมชาติการไม่สลับที่ของการวัดในกลศาสตร์ควอนตัม

พิจารณาการจัดวางเครื่องมือชแตร์น-แกร์ลัคสามชุดติดต่อกัน
\begin{enumerate}
    \item \textbf{เครื่องที่ 1 วางตามแนวแกน $z$ (SGz)} ---- ลำอะตอมแยกออกเป็นสถานะ $\ket{+z}$ และ $\ket{-z}$ เรากั้นลำ $\ket{-z}$ ทิ้ง เหลือเฉพาะลำ $\ket{+z}$ ผ่านเข้าไปยังเครื่องถัดไป
    \item \textbf{เครื่องที่ 2 วางตามแนวแกน $x$ (SGx)} ---- เมื่อลำ $\ket{+z}$ ผ่านสนามแม่เหล็กตามแกน $x$ ลำอนุภาคกลับแยกออกเป็นสองแถบอีกครั้งอย่างเท่าเทียมกัน คือ $\ket{+x}$ (50\%) และ $\ket{-x}$ (50\%)
    \item \textbf{เครื่องที่ 3 วางตามแนวแกน $z$ อีกครั้ง (SGz)} ---- หากเรากั้นลำ $\ket{-x}$ ทิ้ง แล้วนำเฉพาะลำ $\ket{+x}$ ผ่านเข้าสู่ SGz ผลลัพธ์ที่ได้คือ ลำอนุภาคแยกออกเป็น $\ket{+z}$ (50\%) และ $\ket{-z}$ (50\%)
\end{enumerate}

ผลลัพธ์ในขั้นที่ 3 เป็นเรื่องที่ขัดแย้งกับสัญชาตญาณดั้งเดิมอย่างสิ้นเชิง เพราะอนุภาคเคยถูกคัดเลือกให้เป็น $\ket{+z}$ มาแล้วในขั้นที่ 1 แต่การวัดสปินตามแนวแกน $x$ ในขั้นที่ 2 ได้ \textbf{ทำลายข้อมูลสถานะเดิมตามแนวแกน $z$ ไปโดยสิ้นเชิง} เนื่องจาก $[\hat{S}_x, \hat{S}_z] \neq 0$ การวัดในกลศาสตร์ควอนตัมจึงไม่ใช่การสืบค้นคุณสมบัติที่มีอยู่ก่อนอย่างเงียบเชียบ แต่เป็นการรบกวนและบังคับให้ระบบยุบตัวสู่สถานะเจาะจงใหม่ของตัวสังเกตได้นั้น

% ==========================================================
% 6.6 พีชคณิตสปิน เมทริกซ์เพาลี และสปินอร์
% ==========================================================
\section{พีชคณิตสปิน เมทริกซ์เพาลี และสปินอร์}
\index{เมทริกซ์เพาลี (Pauli Spin Matrices)}
\index{สปินอร์ (Spinor)}
\index{ตัวดำเนินการหมุนสปิน (Spin Rotation Operator)}

ในคริสต์ศักราช 1927 โวล์ฟกัง เพาลี (Wolfgang Pauli, 1900--1958) ได้พัฒนาพีชคณิตเมทริกซ์สำหรับระบบควอนตัมสองระดับ (Two-level System) บนปริภูมิฮิลเบิร์ตสองมิติ $\mathcal{H}_2 \cong \mathbb{C}^2$ เพื่ออธิบายสปิน $s = 1/2$

ตัวดำเนินการโมเมนตัมเชิงมุมสปิน $\hat{\mathbf{S}} = (\hat{S}_x, \hat{S}_y, \hat{S}_z)$ สอดคล้องกับพีชคณิตคอมมิวเตเตอร์เช่นเดียวกับโมเมนตัมเชิงมุมทั่วไป
\begin{equation}
[\hat{S}_i, \hat{S}_j] = i\hbar \sum_{k=1}^3 \epsilon_{ijk} \hat{S}_k
\label{eq:spin_commutation}
\end{equation}

เพาลีนิยาม \textbf{เมทริกซ์สปินเพาลี (Pauli Spin Matrices, $\boldsymbol{\sigma}$)} โดยดึงตัวประกอบ $\hbar/2$ ออก
\begin{equation}
\hat{\mathbf{S}} = \frac{\hbar}{2} \boldsymbol{\sigma}
\label{eq:spin_pauli_relation}
\end{equation}
ในฐานเจาะจงของ $\hat{S}_z$ เมทริกซ์เพาลีทั้งสามตัวเขียนได้ในรูปเมทริกซ์มิติ $2 \times 2$
\begin{equation}
\sigma_x = \begin{pmatrix} 0 & 1 \\ 1 & 0 \end{pmatrix}, \qquad \sigma_y = \begin{pmatrix} 0 & -i \\ i & 0 \end{pmatrix}, \qquad \sigma_z = \begin{pmatrix} 1 & 0 \\ 0 & -1 \end{pmatrix}
\label{eq:pauli_matrices_explicit}
\end{equation}

\subsection{เอกลักษณ์และคุณสมบัติพีชคณิตของเมทริกซ์เพาลี}
เมทริกซ์เพาลีมีสมบัติทางพีชคณิตที่งดงามและสำคัญยิ่งยวดในทฤษฎีควอนตัม
\begin{enumerate}
    \item \textbf{ความเป็นแฮร์มิเชียนและเอกภาพ} ---- $\sigma_i^\dagger = \sigma_i$ และ $\sigma_i^2 = I = \begin{pmatrix} 1 & 0 \\ 0 & 1 \end{pmatrix}$ สำหรับทุก $i \in \{x, y, z\}$
    \item \textbf{การแอนติคอมมิวต์ (Anticommutation Relation)}
    \begin{equation}
    \{ \sigma_i, \sigma_j \} \equiv \sigma_i \sigma_j + \sigma_j \sigma_i = 2\delta_{ij} I
    \end{equation}
    \item \textbf{ความสัมพันธ์คอมมิวเตชัน}
    \begin{equation}
    [ \sigma_i, \sigma_j ] \equiv \sigma_i \sigma_j - \sigma_j \sigma_i = 2i \sum_{k=1}^3 \epsilon_{ijk} \sigma_k
    \end{equation}
    \item \textbf{สูตรผลคูณรวมของเพาลี (Master Product Formula)}
    \begin{equation}
    \sigma_i \sigma_j = \delta_{ij} I + i \sum_{k=1}^3 \epsilon_{ijk} \sigma_k
    \label{eq:pauli_master_product}
    \end{equation}
    \item \textbf{ดีเทอร์มิแนนต์และเทรซ} ---- $\det(\sigma_i) = -1$ และ $\text{Tr}(\sigma_i) = 0$ สำหรับทุก $i \in \{x, y, z\}$
\end{enumerate}

จากสูตร~\eqref{eq:pauli_master_product} หากเรามีเวกเตอร์สามมิติทั่วไป $\mathbf{a}$ และ $\mathbf{b}$ จะได้เอกลักษณ์เวกเตอร์ที่ใช้บ่อยที่สุดในฟิสิกส์พลังงานสูงและสปินทรอนิกส์
\begin{equation}
(\boldsymbol{\sigma}\cdot\mathbf{a})(\boldsymbol{\sigma}\cdot\mathbf{b}) = (\mathbf{a}\cdot\mathbf{b}) I + i \boldsymbol{\sigma}\cdot(\mathbf{a}\times\mathbf{b})
\label{eq:pauli_vector_identity}
\end{equation}

\subsection{สปินอร์และการวัดสปินในทิศทางใดๆ}
สถานะสปินทั่วไปของอนุภาคสปินครึ่งหนึ่งเขียนในรูปเวกเตอร์สองมิติเชิงซ้อน เรียกว่า \textbf{สปินอร์ (Spinor, $\chi$)}
\begin{equation}
\chi = \begin{pmatrix} c_+ \\ c_- \end{pmatrix} = c_+ \chi_+ + c_- \chi_- = c_+ \ket{\uparrow} + c_- \ket{\downarrow}
\label{eq:spinor_representation}
\end{equation}
โดยที่ $\chi_+ = \begin{pmatrix} 1 \\ 0 \end{pmatrix}$ คือสถานะสปินขึ้นตามแกน $z$ ($S_z = +\hbar/2$) และ $\chi_- = \begin{pmatrix} 0 \\ 1 \end{pmatrix}$ คือสถานะสปินลงตามแกน $z$ ($S_z = -\hbar/2$) ภายใต้เงื่อนไขนอร์มอลไลซ์ $\chi^\dagger \chi = |c_+|^2 + |c_-|^2 = 1$

หากเราต้องการวัดสปินตามแนวแกนใดๆ ที่กำหนดโดยเวกเตอร์หนึ่งหน่วย $\hat{\mathbf{n}} = (\sin\theta\cos\phi, \sin\theta\sin\phi, \cos\theta)$ ตัวดำเนินการสปินตามแนวแกน $\hat{\mathbf{n}}$ คือ
\begin{equation}
\hat{S}_{\hat{\mathbf{n}}} = \hat{\mathbf{S}}\cdot\hat{\mathbf{n}} = \frac{\hbar}{2} (\boldsymbol{\sigma}\cdot\hat{\mathbf{n}}) = \frac{\hbar}{2} \begin{pmatrix} \cos\theta & \sin\theta e^{-i\phi} \\ \sin\theta e^{i\phi} & -\cos\theta \end{pmatrix}
\label{eq:spin_along_n}
\end{equation}
การแก้สมการค่าเจาะจง $\hat{S}_{\hat{\mathbf{n}}} \chi = \lambda \chi$ ให้ค่าเจาะจง $\lambda = \pm \hbar/2$ โดยมีสถานะเจาะจง (Eigenspinors) คือ
\begin{align}
\chi_+(\hat{\mathbf{n}}) &= \begin{pmatrix} \cos(\theta/2) \\ \sin(\theta/2) e^{i\phi} \end{pmatrix} \label{eq:eigenspinor_plus} \\
\chi_-(\hat{\mathbf{n}}) &= \begin{pmatrix} \sin(\theta/2) \\ -\cos(\theta/2) e^{i\phi} \end{pmatrix} \label{eq:eigenspinor_minus}
\end{align}

\subsection{การหมุนของสปินอร์และความสมมาตรคาบ $4\pi$}
ตัวดำเนินการหมุนสถานะควอนตัมด้วยมุม $\alpha$ รอบแกนหมุน $\hat{\mathbf{n}}$ กำหนดโดยการยกกำลังเอกซ์โพเนนเชียลของตัวกำเนิดการหมุน
\begin{equation}
\hat{U}(\alpha, \hat{\mathbf{n}}) = \exp\left( -\frac{i}{\hbar} \alpha \hat{\mathbf{S}}\cdot\hat{\mathbf{n}} \right) = \exp\left( -i \frac{\alpha}{2} \boldsymbol{\sigma}\cdot\hat{\mathbf{n}} \right)
\end{equation}

โดยการกระจายอนุกรมเทย์เลอร์และใช้สมบัติ $(\boldsymbol{\sigma}\cdot\hat{\mathbf{n}})^2 = I$
\begin{align*}
\exp\left( -i \frac{\alpha}{2} \boldsymbol{\sigma}\cdot\hat{\mathbf{n}} \right) &= \sum_{k=0}^\infty \frac{1}{k!} \left( -i\frac{\alpha}{2} \boldsymbol{\sigma}\cdot\hat{\mathbf{n}} \right)^k \\
&= \sum_{\text{คู่}} \frac{(-1)^{k/2}}{(k)!} \left(\frac{\alpha}{2}\right)^k I - i \sum_{\text{คี่}} \frac{(-1)^{(k-1)/2}}{(k)!} \left(\frac{\alpha}{2}\right)^k (\boldsymbol{\sigma}\cdot\hat{\mathbf{n}})
\end{align*}
จะได้สูตรการหมุนสปินอร์ที่กระชับและสมบูรณ์
\begin{equation}
\hat{U}(\alpha, \hat{\mathbf{n}}) = \cos\left(\frac{\alpha}{2}\right) I - i \sin\left(\frac{\alpha}{2}\right) (\boldsymbol{\sigma}\cdot\hat{\mathbf{n}})
\label{eq:spinor_rotation_operator}
\end{equation}

พิจารณาผลลัพธ์อันน่าพิศวงเมื่อทำการหมุนสปินอร์ไปหนึ่งรอบเต็มในโลกกายภาพ ($\alpha = 2\pi$ หรือ $360^\circ$)
\begin{equation}
\hat{U}(2\pi, \hat{\mathbf{n}}) = \cos(\pi) I - i \sin(\pi) (\boldsymbol{\sigma}\cdot\hat{\mathbf{n}}) = -I
\end{equation}
นั่นคือ \textbf{เมื่อหมุนอนุภาคสปินครึ่งไปครบ $360^\circ$ ฟังก์ชันคลื่นสปินอร์จะไม่กลับมาที่เดิม แต่จะเปลี่ยนเครื่องหมายเป็นลบ ($\chi \to -\chi$)}!

การที่จะทำให้สปินอร์กลับมาสู่สถานะเริ่มต้นอย่างแท้จริง เราต้องหมุนไปถึง \textbf{สองรอบเต็ม ($\alpha = 4\pi$ หรือ $720^\circ$)}
\begin{equation}
\hat{U}(4\pi, \hat{\mathbf{n}}) = \cos(2\pi) I - i \sin(2\pi)(\boldsymbol{\sigma}\cdot\hat{\mathbf{n}}) = +I
\end{equation}
สมบัติคาบ $4\pi$ นี้แสดงว่า การแปลงสถานะสปินครึ่งไม่ได้อยู่ภายใต้กลุ่มการหมุน $SO(3)$ ของเวกเตอร์คลาสสิก แต่อยู่ภายใต้กลุ่มครอบคลุมคู่สองเท่า (Double Covering Group) คือ \textbf{กลุ่มเอกภาพพิเศษ $SU(2)$} ปรากฏการณ์การเปลี่ยนเฟส $\pi$ จากการหมุน $360^\circ$ ได้รับการพิสูจน์ยืนยันความจริงเชิงประจักษ์อย่างหมดจดในการทดลองแทรกสอดลำนิวตรอน (Neutron Interferometry) โดยแอนโทนี แวร์เนอร์ (Anthony Werner) และคณะในปี ค.ศ. 1975

% ==========================================================
% 6.7 สปิน-ออร์บิท โครงสร้างละเอียด และปรากฏการณ์ซีมาน
% ==========================================================
\section{สปิน-ออร์บิท โครงสร้างละเอียด และปรากฏการณ์ซีมาน}
\index{อันตรกิริยาสปิน-ออร์บิท (Spin-Orbit Coupling)}
\index{โครงสร้างละเอียด (Fine Structure)}
\index{ปรากฏการณ์ซีมาน (Zeeman Effect)}

แบบจำลองอะตอมไฮโดรเจนตามสมการชเรอดิงเงอร์ที่ได้ศึกษาในหัวข้อก่อนหน้านี้ ให้ระดับพลังงานที่ขึ้นกับเลขควอนตัม $n$ เพียงอย่างเดียว ทว่าเมื่อวัดสเปกตรัมของไฮโดรเจนด้วยสเปกโตรมิเตอร์ที่มีกำลังแยกขยายสูงยิ่ง จะพบว่าเส้นสเปกตรัมเกิดการแยกตัวออกเป็นเส้นย่อยที่อยู่ชิดกันมาก เรียกว่า \textbf{โครงสร้างละเอียด (Fine Structure)} ซึ่งมีขนาดการแยกประมาณ $\alpha^2 \approx (1/137)^2 \sim 10^{-4}$ ของพลังงานหลัก

โครงสร้างละเอียดเกิดจากผลรวมของการแก้ไขสามประการตามทฤษฎีสัมพัทธภาพพิเศษ
\begin{equation}
\hat{H}_{\text{FS}} = \hat{H}_{\text{rel}} + \hat{H}_{\text{SO}} + \hat{H}_{\text{Darwin}}
\end{equation}

\subsection{การแก้ไขพลังงานจลน์เชิงสัมพัทธภาพ}
ความสัมพันธ์พลังงาน-โมเมนตัมเชิงสัมพัทธภาพคือ $E = \sqrt{p^2 c^2 + m_e^2 c^4}$ เมื่อกระจายทวินามสำหรับ $p \ll m_e c$
\begin{equation}
T = \sqrt{p^2 c^2 + m_e^2 c^4} - m_e c^2 = \frac{p^2}{2m_e} - \frac{p^4}{8m_e^3 c^2} + \dots
\end{equation}
เทอมการรบกวนอันดับหนึ่งของพลังงานจลน์คือ
\begin{equation}
\hat{H}_{\text{rel}} = -\frac{\hat{\mathbf{p}}^4}{8m_e^3 c^2}
\end{equation}

\subsection{อันตรกิริยาสปิน-ออร์บิท (Spin-Orbit Coupling)}
ในกรอบอ้างอิงที่อยู่นิ่งกับอิเล็กตรอน อิเล็กตรอนจะมองเห็นนิวเคลียสโปรตอนเคลื่อนที่วนรอบตัวมันด้วยความเร็ว $-\mathbf{v}$ การเคลื่อนที่ของประจุโปรตอนนี้สร้างกระแสไฟฟ้าและเหนี่ยวนำให้เกิด \textbf{สนามแม่เหล็กภายใน (Internal Magnetic Field, $\mathbf{B}_{\text{int}}$)}
\begin{equation}
\mathbf{B}_{\text{int}} = -\frac{1}{c^2} (\mathbf{v} \times \mathbf{E}) = \frac{1}{m_e c^2 r} \frac{dV}{dr} \mathbf{L}
\end{equation}

โมเมนต์แม่เหล็กสปินของอิเล็กตรอน $\boldsymbol{\mu}_s = -g_s \frac{e}{2m_e}\mathbf{S}$ ย่อมทำอันตรกิริยากับสนามแม่เหล็กภายในนี้ $U = -\boldsymbol{\mu}_s \cdot \mathbf{B}_{\text{int}}$ เมื่อคิดผลของการหมุนควงทอมัส (Thomas Precession) ในทฤษฎีสัมพัทธภาพพิเศษ ซึ่งให้ตัวประกอบลดทอน $1/2$ ฮามิลโทเนียนสปิน-ออร์บิทจึงเขียนได้ดังนี้
\begin{equation}
\hat{H}_{\text{SO}} = \frac{1}{2 m_e^2 c^2} \frac{1}{r} \frac{dV}{dr} (\hat{\mathbf{L}}\cdot\hat{\mathbf{S}})
\label{eq:spin_orbit_hamiltonian}
\end{equation}

สำหรับศักย์คูลอมบ์ $\frac{dV}{dr} = \frac{e^2}{4\pi\varepsilon_0 r^2}$ จะได้
\begin{equation}
\hat{H}_{\text{SO}} = \frac{e^2}{8\pi\varepsilon_0 m_e^2 c^2 r^3} (\hat{\mathbf{L}}\cdot\hat{\mathbf{S}})
\end{equation}

เนื่องจากมีเทอม $\hat{\mathbf{L}}\cdot\hat{\mathbf{S}}$ ตัวดำเนินการ $\hat{L}_z$ และ $\hat{S}_z$ จึงไม่สลับที่กับฮามิลโทเนียนรวมอีกต่อไป เลขควอนตัม $m_\ell$ และ $m_s$ จึงไม่ใช่เลขควอนตัมที่ดี (Good Quantum Numbers) แต่ \textbf{โมเมนตัมเชิงมุมรวม (Total Angular Momentum, $\hat{\mathbf{J}}$)}
\begin{equation}
\hat{\mathbf{J}} \equiv \hat{\mathbf{L}} + \hat{\mathbf{S}}
\end{equation}
จะสลับที่กับฮามิลโทเนียน $[\hat{\mathbf{J}}^2, \hat{H}] = 0$ และ $[\hat{J}_z, \hat{H}] = 0$

เมื่อกระจายขนาดกำลังสอง $\hat{\mathbf{J}}^2 = \hat{\mathbf{L}}^2 + \hat{\mathbf{S}}^2 + 2\hat{\mathbf{L}}\cdot\hat{\mathbf{S}}$ เราสามารถเขียนผลคูณสเกลาร์ได้โดยตรง
\begin{equation}
\hat{\mathbf{L}}\cdot\hat{\mathbf{S}} = \frac{1}{2} \left( \hat{\mathbf{J}}^2 - \hat{\mathbf{L}}^2 - \hat{\mathbf{S}}^2 \right)
\label{eq:l_dot_s_formula}
\end{equation}
ค่าเจาะจงของ $\hat{\mathbf{L}}\cdot\hat{\mathbf{S}}$ ในฐานร่วม $\ket{j, m_j, \ell, s}$ คือ
\begin{equation}
\langle \hat{\mathbf{L}}\cdot\hat{\mathbf{S}} \rangle = \frac{\hbar^2}{2} \big[ j(j + 1) - \ell(\ell + 1) - s(s + 1) \big]
\end{equation}
โดยที่เลขควอนตัมโมเมนตัมเชิงมุมรวม $j$ มีค่าได้สองค่าคือ $j = \ell + 1/2$ หรือ $j = \ell - 1/2$ (สำหรับ $\ell > 0$)

\subsection{สูตรโครงสร้างละเอียดที่สมบูรณ์}
เมื่อรวมผลการแก้ไขเชิงสัมพัทธภาพ การแก้ไขสปิน-ออร์บิท และเทอมดาร์วิน (Darwin Term สำหรับสถานะ $\ell = 0$) ระดับพลังงานโครงสร้างละเอียดของอะตอมไฮโดรเจนเขียนได้ในสูตรที่งดงามอย่างยิ่ง
\begin{equation}
E_{nj} = E_n \left[ 1 + \frac{\alpha^2}{n^2} \left( \frac{n}{j + 1/2} - \frac{3}{4} \right) \right]
\label{eq:fine_structure_full_formula}
\end{equation}
โดยที่ $\alpha \equiv \frac{e^2}{4\pi\varepsilon_0 \hbar c} \approx \frac{1}{137.035999}$ คือ \textbf{ค่าคงตัวโครงสร้างละเอียด (Fine Structure Constant)}

สูตร~\eqref{eq:fine_structure_full_formula} บ่งชี้ว่า ระดับพลังงานที่เคยเสื่อมกันตามค่า $\ell$ ในระดับพลังงานหลัก $n$ เดียวกัน จะเกิดการแยกตัวออกจากกันตามค่าเลขควอนตัมโมเมนตัมเชิงมุมรวม $j$ เช่น ในระดับ $n = 2$ ระดับ $2P_{3/2}$ ($j = 3/2$) จะมีพลังงานสูงกว่าระดับ $2P_{1/2}$ และ $2S_{1/2}$ ($j = 1/2$)

% ==========================================================
% 6.8 การประยุกต์เทคโนโลยีควอนตัมสมัยใหม่ ---- สปินทรอนิกส์และ MRI
% ==========================================================
\section{การประยุกต์เทคโนโลยีควอนตัมสมัยใหม่ ---- สปินทรอนิกส์และ MRI}
\index{สปินทรอนิกส์ (Spintronics)}
\index{ความต้านทานแม่เหล็กยักษ์ (Giant Magnetoresistance)}
\index{การสร้างภาพด้วยเรโซแนนซ์แม่เหล็ก (MRI)}
\index{การหมุนควงลาร์มอร์ (Larmor Precession)}

ความเข้าใจในคุณสมบัติของโมเมนตัมเชิงมุมและสปินของอิเล็กตรอนไม่ได้เป็นเพียงทฤษฎีเชิงนามธรรม แต่เป็นรากฐานที่ขับเคลื่อนเทคโนโลยีเปลี่ยนโลกในศตวรรษที่ 21 ทั้งในด้านการจัดเก็บข้อมูลความหนาแน่นสูง การประมวลผลควอนตัม และการแพทย์วินิจฉัย

\subsection{สปินทรอนิกส์และปรากฏการณ์ความต้านทานแม่เหล็กยักษ์ (GMR)}
ในระบบอิเล็กทรอนิกส์ดั้งเดิม (Conventional Electronics) การประมวลผลและการจัดเก็บสารสนเทศอาศัยการถ่ายโอนประจุไฟฟ้า ($e$) ผ่านสารกึ่งตัวนำ ซึ่งเผชิญกับขีดจำกัดด้านการสูญเสียพลังงานในรูปความร้อนจูล (Joule Heating) และการรั่วไหลของกระแสเมื่ออุปกรณ์มีขนาดเล็กระดับนาโนเมตร

ในทางตรงกันข้าม \textbf{สปินทรอนิกส์ (Spintronics หรือ Spin Electronics)} ใช้ทั้งประจุและสปินของอิเล็กตรอนในการควบคุมการไหลของกระแสไฟฟ้า การค้นพบ \textbf{ปรากฏการณ์ความต้านทานแม่เหล็กยักษ์ (Giant Magnetoresistance, GMR)} ในปี ค.ศ. 1988 โดย อัลแบร์ แฟร์ (Albert Fert, 1938--) และเพเทอร์ กรึนแบร์ก (Peter Grünberg, 1939--2018) ซึ่งได้รับรางวัลโนเบลสาขาฟิสิกส์ในปี ค.ศ. 2007 ได้ปฏิวัติอุตสาหกรรมคอมพิวเตอร์อย่างสิ้นเชิง

โครงสร้าง GMR ประกอบด้วยชั้นโลหะแม่เหล็กเฟอร์โรสองชั้น (เช่น Co, Fe) คั่นกลางด้วยชั้นโลหะที่ไม่เป็นแม่เหล็กความหนาระดับนาโนเมตร (เช่น Cu)
\begin{itemize}
    \item \textbf{เมื่อทิศทางสนามแม่เหล็กของทั้งสองชั้นขนานกัน (Parallel Alignment)} ---- อิเล็กตรอนที่มีทิศทางสปินตรงกับสนามแม่เหล็กจะสามารถเคลื่อนที่ผ่านชั้นรอยต่อได้โดยแทบไม่เกิดการกระเจิง ความต้านทานไฟฟ้าของโครงสร้างจะต่ำมาก ($R_P$)
    \item \textbf{เมื่อทิศทางสนามแม่เหล็กของทั้งสองชั้นตรงข้ามกัน (Antiparallel Alignment)} ---- อิเล็กตรอนทั้งสองทิศทางสปินจะถูกกระเจิงอย่างรุนแรงที่ชั้นรอยต่อชั้นใดชั้นหนึ่งเสมอ ทำให้ความต้านทานไฟฟ้ามีค่าสูงมาก ($R_{AP}$)
\end{itemize}

อัตราส่วนการเปลี่ยนแปลงความต้านทานแม่เหล็กนิยามโดย
\begin{equation}
\text{MR Ratio} = \frac{R_{AP} - R_P}{R_P} \times 100\%
\end{equation}
การค้นพบนี้ถูกนำไปสร้างเป็นหัวอ่านฮาร์ดดิสก์ไดรฟ์ที่ไวต่อสัญญาณแม่เหล็กยิ่งยวด ส่งผลให้ความจุในการจัดเก็บข้อมูลทั่วโลกเพิ่มขึ้นมากกว่าพันเท่า และเป็นต้นกำเนิดของหน่วยความจำเชิงแม่เหล็กเข้าถึงโดยสุ่ม \textbf{MRAM (Magnetoresistive Random Access Memory)} ซึ่งทำงานได้เร็วกว่า SRAM ไม่สูญหายเมื่อตัดกระแสไฟฟ้า (Non-volatile) และทนทานต่อรังสีอวกาศ

\subsection{การสร้างภาพด้วยเรโซแนนซ์แม่เหล็ก (MRI) และการตรวจวัดทางฟิสิกส์การเกษตร}
ในทางการแพทย์และวัสดุศาสตร์ นิวเคลียสของอะตอมที่มีจำนวนโปรตอนหรือนิวตรอนเป็นจำนวนคี่ (โดยเฉพาะอย่างยิ่งนิวเคลียสของอะตอมไฮโดรเจน $^1\text{H}$ ซึ่งเป็นโปรตอนเดี่ยวที่มีอย่างอุดมสมบูรณ์ในน้ำและโมเลกุลชีวภาพ) จะมีโมเมนตัมเชิงมุมสปินนิวเคลียร์ $I = 1/2$

เมื่อวางโปรตอนภายใต้สนามแม่เหล็กหลักสม่ำเสมอ $\mathbf{B}_0 = B_0 \hat{\mathbf{k}}$ โมเมนต์แม่เหล็กนิวเคลียร์ $\boldsymbol{\mu}_N = \gamma \mathbf{I}$ จะแยกออกเป็นสองระดับพลังงานตามปรากฏการณ์ซีมานนิวเคลียร์ โดยมีผลต่างพลังงาน
\begin{equation}
\Delta E = \hbar \gamma B_0
\end{equation}
โดยที่ $\gamma$ คือ \textbf{อัตราส่วนไจโรแมกเนติกของโปรตอน (Gyromagnetic Ratio)} มีค่าเท่ากับ
\begin{equation}
\frac{\gamma}{2\pi} \approx 42.577~\text{MHz/T}
\end{equation}

เวกเตอร์สปินรวมจะหมุนควงรอบทิศทางของสนามแม่เหล็ก $\mathbf{B}_0$ ด้วยความถี่เชิงมุมคงที่ เรียกว่า \textbf{ความถี่ลาร์มอร์ (Larmor Frequency, $\omega_L$)}
\begin{equation}
\omega_L = \gamma B_0
\label{eq:larmor_frequency}
\end{equation}

เมื่อส่งคลื่นแม่เหล็กไฟฟ้าความถี่วิทยุ (Radiofrequency Pulse, RF) ที่มีความถี่ตรงกับ $\omega_L$ เข้าไปในระบบ จะเกิดสภาวะ \textbf{เรโซแนนซ์แม่เหล็กนิวเคลียร์ (Nuclear Magnetic Resonance, NMR)} ดูดกลืนพลังงานและทำให้ทิศทางของสปินสุทธิเบี่ยงเบนออกจากแนวแกน $z$

เมื่อหยุดสัญญาณ RF เวกเตอร์สภาพแม่เหล็กจะคลายตัวกลับสู่จุดสมดุลเดิมตามสมการบล็อค (Bloch Equations) ผ่านสองกลไก
\begin{enumerate}
    \item \textbf{เวลาผ่อนคลายแบบสปิน-แลตทิซ (Spin-Lattice Relaxation Time, $T_1$)} ---- การคืนตัวขององค์ประกอบตามแนวแกนขนาน ($M_z$) สู่สมดุลความร้อน
    \item \textbf{เวลาผ่อนคลายแบบสปิน-สปิน (Spin-Spin Relaxation Time, $T_2$)} ---- การสูญเสียความพร้อมเพรียงของเฟสในระนาบตั้งฉาก ($M_{xy}$)
\end{enumerate}

ในทางการแพทย์ เครื่องตรวจ MRI ในโรงพยาบาล (สนามแม่เหล็ก $1.5\text{ T}$ ถึง $3.0\text{ T}$) อาศัยความแตกต่างของค่า $T_1$ และ $T_2$ ระหว่างเนื้อเยื่อปกติและเนื้อเยื่อที่มีพยาธิสภาพ (เช่น เนื้องอกมะเร็ง สมองขาดเลือด) ในการสร้างภาพสามมิติความละเอียดระดับมิลลิเมตรโดยปราศจากอันตรายจากรังสีเอกซ์

ในด้านฟิสิกส์เกษตรดิจิทัล เทคโนโลยี NMR/MRI ขนาดพกพา (Low-field Portable NMR) ถูกประยุกต์ใช้ในการวัดปริมาณความชื้นในดิน การกระจายตัวของน้ำในท่อน้ำเลี้ยงพืช และการตรวจสอบปริมาณน้ำมันในเมล็ดพันธุ์พืชเศรษฐกิจอย่างรวดเร็วและไม่ทำลายตัวอย่าง ซึ่งเป็นตัวอย่างอันทรงคุณค่าที่แสดงให้เห็นว่า กฎมูลฐานของกลศาสตร์ควอนตัมในระดับอะตอม สามารถเชื่อมโยงไปสู่การพัฒนาเทคโนโลยีการเกษตรแม่นยำเพื่อความมั่นคงทางอาหารได้อย่างเป็นรูปธรรม

% ==========================================================
% ห้องปฏิบัติการเสมือนจริง 3 มิติ WebAR
% ==========================================================
\section{ห้องปฏิบัติการเสมือนจริง 3 มิติ WebAR จำลองออร์บิทัลและชแตร์น-แกร์ลัค}
\index{ห้องปฏิบัติการเสมือนจริง (WebAR Lab)}

เพื่อเสริมสร้างมโนทัศน์เชิงพื้นที่สามมิติ ผู้เรียนสามารถเปิดใช้งานห้องปฏิบัติการเสมือนจริงผ่านเทคโนโลยี Augmented Reality (WebAR) บนสมาร์ทโฟนหรือแท็บเล็ตได้โดยตรงโดยไม่ต้องติดตั้งแอปพลิเคชันเพิ่มเติม

\begin{qrcodelabbox}{ออร์บิทัลอะตอมไฮโดรเจน 3 มิติและการทดลองชแตร์น-แกร์ลัค WebAR}{https://tsanaphysics.github.io/ar-Quantum/lab06_hydrogen_orbitals.html}
\begin{minipage}[c]{0.70\textwidth}
\textbf{ขั้นตอนการทดลองและกิจกรรมการเรียนรู้}
\begin{enumerate}
    \item สแกน QR Code เพื่อเปิดระบบจำลอง 3 มิติบนเว็บเบราว์เซอร์
    \item เลือกแสดงโครงสร้างออร์บิทัล $1s, 2s, 2p_x, 2p_y, 2p_z, 3d_{z^2}, 3d_{x^2-y^2}$ ในรูปกลุ่มหมอกความน่าจะเป็นสามมิติ (Probability Isosurfaces)
    \item ปรับระดับความทึบแสงเพื่อสังเกตบัพเชิงรัศมี (Spherical Shell Nodes) และบัพเชิงมุม (Nodal Planes) ภายในออร์บิทัล
    \item เข้าสู่โหมด ``Stern-Gerlach Beam Splitting Simulator'' ปรับเกรเดียนต์สนามแม่เหล็ก $\partial B_z/\partial z$ และสังเกตการแยกตัวของลำอะตอมเป็นสองทิศทางตามค่า $m_s = \pm 1/2$
    \item ทดลองจัดวางเครื่องมือชแตร์น-แกร์ลัคแบบเรียงลำดับ SGz $\to$ SGx $\to$ SGz เพื่อสังเกตปรากฏการณ์การยุบตัวของสถานะสปินและการทำลายข้อมูลควอนตัม
\end{enumerate}
\end{minipage}
\hfill
\begin{minipage}[c]{0.25\textwidth}
\centering
\qrcode[height=2.8cm]{https://tsanaphysics.github.io/ar-Quantum/lab06_hydrogen_orbitals.html}
\end{minipage}
\end{qrcodelabbox}

% ==========================================================
% ตัวอย่างการแก้โจทย์เชิงวิเคราะห์ตามระเบียบวิธี P-S-C-T
% ==========================================================
\section{ตัวอย่างการแก้โจทย์เชิงวิเคราะห์ตามระเบียบวิธี P-S-C-T}
\index{ระเบียบวิธี P-S-C-T (P-S-C-T Method)}

การทำความเข้าใจฟิสิกส์ควอนตัมอย่างลึกซึ้งจำเป็นต้องอาศัยการฝึกฝนทักษะการแก้ปัญหาเชิงปริมาณอย่างเป็นระบบ ระเบียบวิธี \textbf{P-S-C-T} ประกอบด้วย 4 ขั้นตอนหลัก
\begin{itemize}
    \item \textbf{P (Problem)} ---- ระบุแก่นแท้ของโจทย์ เงื่อนไขขอบเขต และเป้าหมายทางฟิสิกส์
    \item \textbf{S (Solution)} ---- แสดงขั้นตอนการอนุพัทธ์ทางคณิตศาสตร์อย่างโปร่งใส เป็นลำดับขั้น
    \item \textbf{C (Check)} ---- ตรวจสอบความถูกต้องของหน่วย มิติ พฤติกรรมที่จุดลิมิตสุดขั้ว
    \item \textbf{T (Think)} ---- วิเคราะห์ความหมายทางฟิสิกส์ การเชื่อมโยงสู่ปรากฏการณ์ระดับมหภาค
\end{itemize}

% ----------------------------------------------------------
% ตัวอย่างที่ 6.1
% ----------------------------------------------------------
\begin{psctbox}{ตัวอย่างที่ 6.1 ---- ความน่าจะเป็นในการค้นพบอิเล็กตรอนภายในทรงกลมรัศมีโบร์}

\textbf{Problem (โจทย์ปัญหา)}

สำหรับอิเล็กตรอนในสถานะพื้น ($1s, n=1, \ell=0, m=0$) ของอะตอมไฮโดรเจน ฟังก์ชันคลื่นนอร์มอลไลซ์คือ
\begin{equation*}
\psi_{100}(r) = \frac{1}{\sqrt{\pi a_0^3}} e^{-r/a_0}
\end{equation*}
\begin{enumerate}
    \item จงหาระยะรัศมี $r_{\text{max}}$ ที่มีโอกาสพบอิเล็กตรอนมากที่สุด (Most Probable Radius)
    \item จงคำนวณความน่าจะเป็นที่อิเล็กตรอนจะอยู่ภายในทรงกลมรัศมีโบร์ ($r \le a_0$)
\end{enumerate}

\textbf{Solution (วิธีทำ)}

\textbf{ขั้นตอนที่ 1} หาความหนาแน่นความน่าจะเป็นในแนวรัศมี $P(r)$
\begin{equation*}
P(r) = r^2 |R_{10}(r)|^2 = 4\pi r^2 |\psi_{100}|^2 = \frac{4}{a_0^3} r^2 e^{-2r/a_0}
\end{equation*}

\textbf{ขั้นตอนที่ 2} หาระยะ $r_{\text{max}}$ จากเงื่อนไขค่าสูงสุดสัมพัทธ์ $\frac{dP}{dr} = 0$
\begin{equation*}
\frac{dP}{dr} = \frac{4}{a_0^3} \left[ 2r e^{-2r/a_0} + r^2 \left(-\frac{2}{a_0}\right) e^{-2r/a_0} \right] = \frac{8}{a_0^3} r \left( 1 - \frac{r}{a_0} \right) e^{-2r/a_0} = 0
\end{equation*}
เนื่องจาก $r \neq 0$ และ $e^{-2r/a_0} \neq 0$ จะได้
\begin{equation*}
1 - \frac{r}{a_0} = 0 \quad \implies \quad r_{\text{max}} = a_0
\end{equation*}
ระยะทางที่มีความหนาแน่นความน่าจะเป็นสูงสุดมีค่าเท่ากับรัศมีโบร์ $a_0$ พอดี

\textbf{ขั้นตอนที่ 3} คำนวณความน่าจะเป็น $P(r \le a_0)$ โดยการอินทิเกรต
\begin{equation*}
P(r \le a_0) = \int_0^{a_0} P(r) dr = \frac{4}{a_0^3} \int_0^{a_0} r^2 e^{-2r/a_0} dr
\end{equation*}
เปลี่ยนตัวแปร $x = 2r/a_0$ ดังนั้น $r = \frac{a_0}{2} x$ และ $dr = \frac{a_0}{2} dx$ โดยขอบเขตเปลี่ยนเป็น $x = 0$ ถึง $x = 2$
\begin{equation*}
P(r \le a_0) = \frac{4}{a_0^3} \int_0^2 \left(\frac{a_0^2}{4} x^2\right) e^{-x} \left(\frac{a_0}{2} dx\right) = \frac{1}{2} \int_0^2 x^2 e^{-x} dx
\end{equation*}
ใช้อินทิเกรตทีละส่วน (Integration by Parts) หรือสูตรมาตรฐาน $\int x^2 e^{-x} dx = -e^{-x}(x^2 + 2x + 2)$
\begin{align*}
P(r \le a_0) &= \frac{1}{2} \Big[ -e^{-x}(x^2 + 2x + 2) \Big]_0^2 \\
&= \frac{1}{2} \left[ -e^{-2}(4 + 4 + 2) - (-e^0(2)) \right] \\
&= \frac{1}{2} \left[ 2 - 10 e^{-2} \right] = 1 - 5 e^{-2}
\end{align*}
แทนค่า $e^{-2} \approx 0.135335$
\begin{equation*}
P(r \le a_0) = 1 - 5(0.135335) = 1 - 0.676676 \approx 0.3233 = 32.33\%
\end{equation*}

\textbf{Check (ตรวจสอบ)} ---- ค่าความน่าจะเป็นอยู่ระหว่าง 0 ถึง 1 ($0.3233 \in [0, 1]$) ซึ่งถูกต้องเชิงความน่าจะเป็น และมีค่าน้อยกว่า 50\% สอดคล้องกับธรรมชาติของการแจกแจงที่เบ้ขวา

\textbf{Think (คิดต่อ)} ---- แม้ว่าแบบจำลองอะตอมของบอร์จะระบุว่าอิเล็กตรอนต้องโคจรอยู่ที่รัศมี $a_0$ อย่างแน่นอน แต่ในกลศาสตร์ควอนตัม โอกาสที่จะพบอิเล็กตรอนอยู่ด้านในทรงกลมรัศมีโบร์มีเพียง $32.3\%$ เท่านั้น ส่วนอีก $67.7\%$ อิเล็กตรอนจะกระจายตัวอยู่ภายนอกรัศมีโบร์
\end{psctbox}

% ----------------------------------------------------------
% ตัวอย่างที่ 6.2
% ----------------------------------------------------------
\begin{psctbox}{ตัวอย่างที่ 6.2 ---- บัพเชิงรัศมีและค่าคาดหวังของระยะทางสำหรับสถานะ $2s$ และ $2p$}

\textbf{Problem (โจทย์ปัญหา)}

ฟังก์ชันคลื่นเชิงรัศมีของสถานะ $2s$ และ $2p$ ของอะตอมไฮโดรเจนคือ
\begin{equation*}
R_{20}(r) = \frac{1}{2\sqrt{2}} a_0^{-3/2} \left( 2 - \frac{r}{a_0} \right) e^{-r/2a_0}, \qquad R_{21}(r) = \frac{1}{2\sqrt{6}} a_0^{-3/2} \left( \frac{r}{a_0} \right) e^{-r/2a_0}
\end{equation*}
\begin{enumerate}
    \item จงหาตำแหน่งบัพเชิงรัศมีของสถานะ $2s$
    \item จงคำนวณค่าคาดหวังของระยะทาง $\langle r \rangle_{2s}$ และ $\langle r \rangle_{2p}$ พร้อมเปรียบเทียบผลลัพธ์
\end{enumerate}

\textbf{Solution (วิธีทำ)}

\textbf{ขั้นตอนที่ 1} หาตำแหน่งบัพเชิงรัศมีของสถานะ $2s$
บัพเชิงรัศมีเกิดขึ้นเมื่อ $R_{20}(r) = 0$
\begin{equation*}
2 - \frac{r}{a_0} = 0 \quad \implies \quad r_{\text{node}} = 2 a_0
\end{equation*}
ดังนั้น ความน่าจะเป็นในการพบอิเล็กตรอนเป็นศูนย์บนผิวเปลือกทรงกลมรัศมี $2a_0$

\textbf{ขั้นตอนที่ 2} คำนวณค่าคาดหวัง $\langle r \rangle_{2s}$
\begin{equation*}
\langle r \rangle_{2s} = \int_0^\infty r P_{20}(r) dr = \int_0^\infty r^3 |R_{20}(r)|^2 dr = \frac{1}{8 a_0^3} \int_0^\infty r^3 \left( 2 - \frac{r}{a_0} \right)^2 e^{-r/a_0} dr
\end{equation*}
กระจายเทอมในวงเล็บ $\left( 2 - \frac{r}{a_0} \right)^2 = 4 - \frac{4r}{a_0} + \frac{r^2}{a_0^2}$
\begin{equation*}
\langle r \rangle_{2s} = \frac{1}{8 a_0^3} \left[ 4 \int_0^\infty r^3 e^{-r/a_0} dr - \frac{4}{a_0} \int_0^\infty r^4 e^{-r/a_0} dr + \frac{1}{a_0^2} \int_0^\infty r^5 e^{-r/a_0} dr \right]
\end{equation*}
ใช้อินทิกรัลมาตรฐาน $\int_0^\infty r^m e^{-\alpha r} dr = \frac{m!}{\alpha^{m+1}}$ โดย $\alpha = 1/a_0$
\begin{align*}
\langle r \rangle_{2s} &= \frac{1}{8 a_0^3} \left[ 4(6 a_0^4) - \frac{4}{a_0}(24 a_0^5) + \frac{1}{a_0^2}(120 a_0^6) \right] \\
&= \frac{a_0}{8} \big[ 24 - 96 + 120 \big] = \frac{a_0}{8} (48) = 6 a_0
\end{align*}

\textbf{ขั้นตอนที่ 3} คำนวณค่าคาดหวัง $\langle r \rangle_{2p}$
\begin{equation*}
\langle r \rangle_{2p} = \int_0^\infty r^3 |R_{21}(r)|^2 dr = \frac{1}{24 a_0^5} \int_0^\infty r^5 e^{-r/a_0} dr = \frac{1}{24 a_0^5} (120 a_0^6) = 5 a_0
\end{equation*}

\textbf{Check (ตรวจสอบ)} ---- ตรวจสอบกับสูตรทั่วไปของคราเมอส์ สมการ~\eqref{eq:exp_r} $\langle r \rangle_{n\ell} = \frac{a_0}{2}[3n^2 - \ell(\ell+1)]$
\begin{itemize}
    \item สำหรับ $2s$ ($n=2, \ell=0$) ---- $\langle r \rangle = \frac{a_0}{2}[3(4) - 0] = 6 a_0$ ตรงกันแม่นยำ
    \item สำหรับ $2p$ ($n=2, \ell=1$) ---- $\langle r \rangle = \frac{a_0}{2}[3(4) - 2] = 5 a_0$ ตรงกันแม่นยำ
\end{itemize}

\textbf{Think (คิดต่อ)} ---- แม้ว่า $\langle r \rangle_{2s} = 6 a_0 > \langle r \rangle_{2p} = 5 a_0$ แต่อิเล็กตรอน $2s$ มีพูความหนาแน่นชั้นในใกล้จุดกำเนิด ($r < 2a_0$) ทำให้มีโอกาสแทรกซึม (Penetration) เข้าใกล้นิวเคลียสได้มากกว่าอิเล็กตรอน $2p$ ปรากฏการณ์นี้มีความสำคัญยิ่งยวดในอะตอมหลายอิเล็กตรอน เพราะทำให้อิเล็กตรอน $2s$ ถูกกำบังประจุบวกน้อยกว่าและมีพลังงานต่ำกว่า $2p$
\end{psctbox}

% ----------------------------------------------------------
% ตัวอย่างที่ 6.3
% ----------------------------------------------------------
\begin{psctbox}{ตัวอย่างที่ 6.3 ---- การวัดสปินแบบเรียงลำดับและการยุบตัวของสถานะควอนตัม}

\textbf{Problem (โจทย์ปัญหา)}

ลำอะตอมเงินที่มีสปิน $1/2$ ถูกส่งผ่านระบบเครื่องมือชแตร์น-แกร์ลัคสามชุดติดต่อกัน
\begin{enumerate}
    \item เครื่องที่หนึ่ง ($SG_z$) ---- คัดเลือกเฉพาะสถานะสปินขึ้นตามแนวแกน $z$ ($\ket{\uparrow_z}$) ออกมา
    \item เครื่องที่สอง ($SG_\theta$) ---- วางทำมุม $\theta$ เทียบกับแกน $z$ ในระนาบ $xz$ โดยคัดเลือกเฉพาะสถานะสปินขึ้น $\ket{\uparrow_\theta}$ ออกมา
    \item เครื่องที่สาม ($SG_z$) ---- วัดสปินตามแนวแกน $z$ อีกครั้งหนึ่ง
\end{enumerate}
จงคำนวณสัดส่วนของอะตอมทั้งหมดที่หลุดรอดผ่านระบบออกมาได้ในสถานะสปินลงตามแนวแกน $z$ ($\ket{\downarrow_z}$) หลังผ่านเครื่องที่สาม

\textbf{Solution (วิธีทำ)}

\textbf{ขั้นตอนที่ 1} เขียนสถานะเจาะจงของสปินตามแนวแกน $\theta$
แกนการวัดชี้ตามเวกเตอร์ $\hat{\mathbf{n}} = (\sin\theta, 0, \cos\theta)$ ซึ่งเทียบเท่ากับ $\phi = 0$ ในสมการ~\eqref{eq:eigenspinor_plus}
\begin{equation*}
\ket{\uparrow_\theta} = \cos(\theta/2) \ket{\uparrow_z} + \sin(\theta/2) \ket{\downarrow_z}
\end{equation*}

\textbf{ขั้นตอนที่ 2} คำนวณความน่าจะเป็นที่จะผ่านเครื่องที่สอง ($SG_\theta$)
อะตอมที่เข้าสู่เครื่องที่สองอยู่ในสถานะ $\ket{\uparrow_z}$ แอมพลิจูดความน่าจะเป็นในการวัดได้สถานะ $\ket{\uparrow_\theta}$ คือ
\begin{equation*}
\braket{\uparrow_\theta}{\uparrow_z} = \cos(\theta/2)
\end{equation*}
ความน่าจะเป็นที่จะผ่านเครื่องที่สองคือ
\begin{equation*}
P_2 = |\braket{\uparrow_\theta}{\uparrow_z}|^2 = \cos^2(\theta/2)
\end{equation*}

\textbf{ขั้นตอนที่ 3} คำนวณความน่าจะเป็นที่จะผ่านเครื่องที่สามและได้สถานะ $\ket{\downarrow_z}$
อะตอมที่หลุดออกจากเครื่องที่สองจะยุบตัวลงสู่สถานะ $\ket{\uparrow_\theta}$ โดยสมบูรณ์ แอมพลิจูดความน่าจะเป็นในการวัดได้สถานะ $\ket{\downarrow_z}$ คือ
\begin{equation*}
\braket{\downarrow_z}{\uparrow_\theta} = \sin(\theta/2)
\end{equation*}
ความน่าจะเป็นคือ
\begin{equation*}
P_3 = |\braket{\downarrow_z}{\uparrow_\theta}|^2 = \sin^2(\theta/2)
\end{equation*}

\textbf{ขั้นตอนที่ 4} คำนวณความน่าจะเป็นรวมตลอดทั้งระบบ
\begin{equation*}
P_{\text{total}} = P_2 \times P_3 = \cos^2(\theta/2) \sin^2(\theta/2) = \frac{1}{4} \sin^2\theta
\end{equation*}

\textbf{Check (ตรวจสอบ)}
\begin{itemize}
    \item ถ้า $\theta = 0$ (เครื่องที่สองวางตามแกน $z$ เช่นเดิม) ---- $P_{\text{total}} = \frac{1}{4}\sin^2(0) = 0$ สอดคล้องกับความเป็นจริง เพราะสถานะ $\ket{\uparrow_z}$ จะไม่มีทางวัดได้ $\ket{\downarrow_z}$
    \item ถ้า $\theta = \pi/2$ (เครื่องที่สองวางตามแกน $x$) ---- $P_{\text{total}} = \frac{1}{4}\sin^2(\pi/2) = \frac{1}{4} = 25\%$ ซึ่งตรงกับผลคูณ $50\% \times 50\% = 25\%$
\end{itemize}

\textbf{Think (คิดต่อ)} ---- การวางเครื่องมือที่สองคั่นกลางทำให้ระบบสามารถเกิดสถานะ $\ket{\downarrow_z}$ ขึ้นมาได้ ทั้งๆ ที่ในตอนแรกอะตอมทั้งหมดเป็น $\ket{\uparrow_z}$ อย่างแน่นอน ค่าความน่าจะเป็นสูงสุดจะเกิดขึ้นเมื่อ $\theta = \pi/2$ ซึ่งเป็นหลักการพื้นฐานที่ใช้ในการสื่อสารควอนตัมและการเข้ารหัสลับ BB84
\end{psctbox}

% ----------------------------------------------------------
% ตัวอย่างที่ 6.4
% ----------------------------------------------------------
\begin{psctbox}{ตัวอย่างที่ 6.4 ---- การหมุนสปินอร์และการเปลี่ยนเครื่องหมายภายใต้การหมุน $360^\circ$}

\textbf{Problem (โจทย์ปัญหา)}

กำหนดให้อนุภาคสปินครึ่งตัวหนึ่งเริ่มต้นอยู่ในสถานะสปินขึ้นตามแนวแกน $x$ กล่าวคือ $\chi(0) = \ket{\uparrow_x} = \frac{1}{\sqrt{2}}\begin{pmatrix} 1 \\ 1 \end{pmatrix}$
\begin{enumerate}
    \item หากทำการหมุนสถานะนี้รอบแกน $z$ เป็นมุม $\alpha$ จงหาสปินอร์ลัพธ์ $\chi(\alpha)$
    \item จงแสดงว่าเมื่อหมุนไปเป็นมุม $\alpha = 2\pi$ ($360^\circ$) สถานะสปินอร์จะเปลี่ยนเครื่องหมายเป็นลบ และคำนวณค่าคาดหวัง $\langle S_x \rangle$ ที่มุม $\alpha = \pi/2$
\end{enumerate}

\textbf{Solution (วิธีทำ)}

\textbf{ขั้นตอนที่ 1} ใช้ตัวดำเนินการหมุนสปินอร์รอบแกน $z$
จากสมการ~\eqref{eq:spinor_rotation_operator} โดยกำหนด $\hat{\mathbf{n}} = \hat{\mathbf{k}}$
\begin{equation*}
\hat{U}(\alpha, \hat{\mathbf{k}}) = \cos(\alpha/2) I - i \sin(\alpha/2) \sigma_z = \begin{pmatrix} e^{-i\alpha/2} & 0 \\ 0 & e^{i\alpha/2} \end{pmatrix}
\end{equation*}

\textbf{ขั้นตอนที่ 2} หาสปินอร์ลัพธ์ $\chi(\alpha)$
\begin{equation*}
\chi(\alpha) = \hat{U}(\alpha, \hat{\mathbf{k}}) \chi(0) = \begin{pmatrix} e^{-i\alpha/2} & 0 \\ 0 & e^{i\alpha/2} \end{pmatrix} \frac{1}{\sqrt{2}} \begin{pmatrix} 1 \\ 1 \end{pmatrix} = \frac{1}{\sqrt{2}} \begin{pmatrix} e^{-i\alpha/2} \\ e^{i\alpha/2} \end{pmatrix}
\end{equation*}

\textbf{ขั้นตอนที่ 3} ตรวจสอบการหมุนหนึ่งรอบเต็ม ($\alpha = 2\pi$)
\begin{equation*}
\chi(2\pi) = \frac{1}{\sqrt{2}} \begin{pmatrix} e^{-i\pi} \\ e^{i\pi} \end{pmatrix} = \frac{1}{\sqrt{2}} \begin{pmatrix} -1 \\ -1 \end{pmatrix} = -\chi(0)
\end{equation*}
สปินอร์เปลี่ยนเครื่องหมายเป็นลบอย่างแท้จริง

\textbf{ขั้นตอนที่ 4} คำนวณค่าคาดหวัง $\langle S_x \rangle$ ที่มุม $\alpha = \pi/2$
ที่ $\alpha = \pi/2$ จะได้ $\chi(\pi/2) = \frac{1}{\sqrt{2}}\begin{pmatrix} e^{-i\pi/4} \\ e^{i\pi/4} \end{pmatrix}$
\begin{align*}
\langle S_x \rangle &= \chi^\dagger \left( \frac{\hbar}{2}\sigma_x \right) \chi = \frac{\hbar}{4} \begin{pmatrix} e^{i\pi/4} & e^{-i\pi/4} \end{pmatrix} \begin{pmatrix} 0 & 1 \\ 1 & 0 \end{pmatrix} \begin{pmatrix} e^{-i\pi/4} \\ e^{i\pi/4} \end{pmatrix} \\
&= \frac{\hbar}{4} \begin{pmatrix} e^{i\pi/4} & e^{-i\pi/4} \end{pmatrix} \begin{pmatrix} e^{i\pi/4} \\ e^{-i\pi/4} \end{pmatrix} = \frac{\hbar}{4} \big[ e^{i\pi/2} + e^{-i\pi/2} \big] \\
&= \frac{\hbar}{4} \big[ 2 \cos(\pi/2) \big] = 0
\end{align*}

\textbf{Check (ตรวจสอบ)} ---- เมื่อเริ่มต้นสปินชี้ตามแนวแกน $+x$ แล้วหมุนรอบแกน $z$ ไป $90^\circ$ เวกเตอร์สปินจะชี้ไปตามแนวแกน $+y$ ทำให้องค์ประกอบตามแนวแกน $x$ ต้องมีค่าเป็นศูนย์พอดี ซึ่งตรงกับ $\langle S_x \rangle = 0$

\textbf{Think (คิดต่อ)} ---- แม้ว่าการเปลี่ยนเครื่องหมาย $\chi \to -\chi$ จะไม่เปลี่ยนค่าคาดหวังของตัวดำเนินการเดี่ยว (เนื่องจาก $\bra{-\chi}\hat{O}\ket{-\chi} = \bra{\chi}\hat{O}\ket{\chi}$) แต่หากนำลำนิวตรอนที่ถูกหมุน $360^\circ$ ไปแทรกสอดกับลำอ้างอิงที่ไม่ถูกหมุน เฟส $-\pi$ จะทำให้เกิดการแทรกสอดแบบหักล้างอย่างสมบูรณ์ ซึ่งเป็นการทดลองจริงที่พิสูจน์ธรรมชาติควอนตัมของสปิน
\end{psctbox}

% ----------------------------------------------------------
% ตัวอย่างที่ 6.5
% ----------------------------------------------------------
\begin{psctbox}{ตัวอย่างที่ 6.5 ---- การแยกตัวของระดับพลังงานจากอันตรกิริยาสปิน-ออร์บิทในระดับ $2p$}

\textbf{Problem (โจทย์ปัญหา)}

สำหรับระดับพลังงาน $n = 2, \ell = 1$ ($2p$) ของอะตอมไฮโดรเจน อิเล็กตรอนมีสปิน $s = 1/2$
\begin{enumerate}
    \item จงหาค่าของโมเมนตัมเชิงมุมรวม $j$ ที่เป็นไปได้ทั้งหมด และระบุสัญลักษณ์สเปกโทรสโกปี $^{2s+1}L_j$
    \item จงคำนวณผลต่างพลังงานสปิน-ออร์บิท $\Delta E_{\text{SO}}$ ระหว่างสองระดับนี้ในหน่วยอิเล็กตรอนโวลต์ ($\text{eV}$)
\end{enumerate}

\textbf{Solution (วิธีทำ)}

\textbf{ขั้นตอนที่ 1} หาค่า $j$ และสัญลักษณ์สเปกโทรสโกปี
จากกฎการบวกโมเมนตัมเชิงมุม $j = |\ell - s|, \dots, \ell + s$
\begin{equation*}
j = 1 - 1/2 = 1/2 \qquad \text{และ} \qquad j = 1 + 1/2 = 3/2
\end{equation*}
ระดับ $2p$ จึงแยกออกเป็นสองระดับย่อย คือ \textbf{$^2P_{1/2}$} ($j = 1/2$) และ \textbf{$^2P_{3/2}$} ($j = 3/2$)

\textbf{ขั้นตอนที่ 2} คำนวณค่าคาดหวังของ $\hat{\mathbf{L}}\cdot\hat{\mathbf{S}}$ สำหรับแต่ละระดับ
จากสมการ~\eqref{eq:l_dot_s_formula}
\begin{itemize}
    \item สำหรับ $^2P_{3/2}$ ($j = 3/2, \ell = 1, s = 1/2$)
    \begin{equation*}
    \langle \hat{\mathbf{L}}\cdot\hat{\mathbf{S}} \rangle = \frac{\hbar^2}{2} \left[ \frac{3}{2}\left(\frac{5}{2}\right) - 1(2) - \frac{1}{2}\left(\frac{3}{2}\right) \right] = \frac{\hbar^2}{2} \left[ \frac{15}{4} - 2 - \frac{3}{4} \right] = \frac{\hbar^2}{2} (1) = \frac{\hbar^2}{2}
    \end{equation*}
    \item สำหรับ $^2P_{1/2}$ ($j = 1/2, \ell = 1, s = 1/2$)
    \begin{equation*}
    \langle \hat{\mathbf{L}}\cdot\hat{\mathbf{S}} \rangle = \frac{\hbar^2}{2} \left[ \frac{1}{2}\left(\frac{3}{2}\right) - 2 - \frac{3}{4} \right] = \frac{\hbar^2}{2} \left[ \frac{3}{4} - 2 - \frac{3}{4} \right] = -\hbar^2
    \end{equation*}
\end{itemize}

\textbf{ขั้นตอนที่ 3} คำนวณพลังงานสปิน-ออร์บิท
\begin{equation*}
E_{\text{SO}} = \frac{e^2}{8\pi\varepsilon_0 m_e^2 c^2} \left\langle \frac{1}{r^3} \right\rangle_{21} \langle \hat{\mathbf{L}}\cdot\hat{\mathbf{S}} \rangle
\end{equation*}
จากสมการ~\eqref{eq:exp_inv_r3} $\langle 1/r^3 \rangle_{21} = \frac{1}{2^3 (1)(3/2)(2) a_0^3} = \frac{1}{24 a_0^3}$
ผลต่างพลังงานระหว่างระดับ $^2P_{3/2}$ และ $^2P_{1/2}$ คือ
\begin{equation*}
\Delta E_{\text{SO}} = E_{\text{SO}}(3/2) - E_{\text{SO}}(1/2) = \frac{e^2 \hbar^2}{8\pi\varepsilon_0 m_e^2 c^2 (24 a_0^3)} \left[ \frac{1}{2} - (-1) \right] = \frac{3}{2} \frac{e^2 \hbar^2}{192 \pi\varepsilon_0 m_e^2 c^2 a_0^3}
\end{equation*}
เขียนในรูปค่าคงตัวโครงสร้างละเอียด $\alpha = \frac{e^2}{4\pi\varepsilon_0 \hbar c}$ และพลังงานสถานะพื้น $|E_1| = \frac{\hbar^2}{2m_e a_0^2} \approx 13.6~\text{eV}$
\begin{equation*}
\Delta E_{\text{SO}} = \frac{\alpha^2 |E_1|}{16} \approx \frac{(1/137.036)^2 (13.606~\text{eV})}{16} \approx 4.53 \times 10^{-5}~\text{eV}
\end{equation*}

\textbf{Check (ตรวจสอบ)} ---- มิติของผลลัพธ์เป็นหน่วยพลังงาน และมีอันดับขนาดสอดคล้องกับโครงสร้างละเอียด ($\sim 10^{-5}~\text{eV}$) ซึ่งตรงกับการสังเกตการณ์เส้นสเปกตรัมไฮโดรเจนอย่างแม่นยำ

\textbf{Think (คิดต่อ)} ---- การแยกตัวของระดับพลังงานเพียง $4.53 \times 10^{-5}~\text{eV}$ นี้ แม้จะดูเล็กน้อยเมื่อเทียบกับพลังงานหลัก $13.6~\text{eV}$ แต่เป็นตัวการสำคัญที่ทำให้เกิดเส้นคู่ (Doublet) ในสเปกตรัมของโลหะแอลคาไล เช่น เส้นสีเหลืองคู่โซเดียม ($D_1$ และ $D_2$) ที่เราพบเห็นในหลอดไฟตามท้องถนน
\end{psctbox}

% ----------------------------------------------------------
% ตัวอย่างที่ 6.6
% ----------------------------------------------------------
\begin{psctbox}{ตัวอย่างที่ 6.6 ---- สนามแม่เหล็กภายในที่อิเล็กตรอนสัมผัสในสถานะ $2p$}

\textbf{Problem (โจทย์ปัญหา)}

จากผลต่างพลังงานสปิน-ออร์บิทในตัวอย่างที่ 6.5 จงคำนวณขนาดของสนามแม่เหล็กยังผลภายใน ($\mathbf{B}_{\text{int}}$) ที่อิเล็กตรอนสัมผัสเนื่องจากการโคจรสัมพัทธ์รอบโปรตอนในสถานะ $2p$

\textbf{Solution (วิธีทำ)}

\textbf{ขั้นตอนที่ 1} ตั้งสมการพลังงานขั้วคู่แม่เหล็ก
พลังงานปฏิสัมพันธ์ระหว่างสปินอิเล็กตรอนและสนามแม่เหล็กภายในคือ
\begin{equation*}
U = -\boldsymbol{\mu}_s \cdot \mathbf{B}_{\text{int}} = g_s \frac{\mu_B}{\hbar} \mathbf{S}\cdot\mathbf{B}_{\text{int}} \approx 2 \frac{\mu_B}{\hbar} S_z B_{\text{int}} = \pm \mu_B B_{\text{int}}
\end{equation*}
ผลต่างพลังงานระหว่างสปินขนานและตรงข้ามสนามแม่เหล็กคือ
\begin{equation*}
\Delta E_{\text{SO}} = 2 \mu_B B_{\text{int}}
\end{equation*}

\textbf{ขั้นตอนที่ 2} คำนวณขนาดของสนามแม่เหล็ก $B_{\text{int}}$
\begin{equation*}
B_{\text{int}} = \frac{\Delta E_{\text{SO}}}{2 \mu_B}
\end{equation*}
แทนค่า $\Delta E_{\text{SO}} = 4.53 \times 10^{-5}~\text{eV} = 4.53 \times 10^{-5} \times 1.602 \times 10^{-19}~\text{J} \approx 7.257 \times 10^{-24}~\text{J}$ และโบร์แมกเนตอน $\mu_B = 9.274 \times 10^{-24}~\text{J/T}$
\begin{equation*}
B_{\text{int}} = \frac{7.257 \times 10^{-24}~\text{J}}{2(9.274 \times 10^{-24}~\text{J/T})} \approx 0.391~\text{Tesla}
\end{equation*}

\textbf{Check (ตรวจสอบ)} ---- สนามแม่เหล็กมีหน่วยเป็นเทสลา (Tesla) และมีขนาดประมาณ $0.4\text{ T}$ ซึ่งเป็นขนาดสนามแม่เหล็กในห้องปฏิบัติการทั่วไป

\textbf{Think (คิดต่อ)} ---- สนามแม่เหล็กภายในอะตอมขนาดประมาณ $0.4\text{ T}$ นี้ มีความแรงมากกว่าสนามแม่เหล็กโลก ($0.00005\text{ T}$) ถึงเกือบหนึ่งหมื่นเท่า! แสดงให้เห็นว่าแม้ในอะตอมที่เรียบง่ายอย่างไฮโดรเจน อิเล็กตรอนเคลื่อนที่ด้วยความเร็วสูงมากในบริเวณขนาดเล็ก จนเหนี่ยวนำให้เกิดสนามแม่เหล็กมหาศาลภายในตัวเอง
\end{psctbox}

% ----------------------------------------------------------
% ตัวอย่างที่ 6.7
% ----------------------------------------------------------
\begin{psctbox}{ตัวอย่างที่ 6.7 ---- ความถี่ลาร์มอร์และการเรโซแนนซ์สปินในเครื่องตรวจ MRI}

\textbf{Problem (โจทย์ปัญหา)}

เครื่องตรวจวินิจฉัยทางการแพทย์ MRI รุ่นมาตรฐานในโรงพยาบาล ใช้สนามแม่เหล็กตัวนำยิ่งยวดขนาด $B_0 = 3.0~\text{Tesla}$
\begin{enumerate}
    \item จงคำนวณความถี่ลาร์มอร์ ($\nu_L = \omega_L/2\pi$) ของโปรตอนในร่างกายมนุษย์
    \item หากความแปรปรวนของสนามแม่เหล็กในบริเวณตรวจวัดคือ $\Delta B = 1.0 \times 10^{-4}~\text{Tesla}$ ($1\text{ Gauss}$) จงหาความกว้างของแถบความถี่เรโซแนนซ์ $\Delta \nu_L$ ที่ระบบต้องรองรับ
\end{enumerate}

\textbf{Solution (วิธีทำ)}

\textbf{ขั้นตอนที่ 1} คำนวณความถี่ลาร์มอร์ของโปรตอน
จากสมการ~\eqref{eq:larmor_frequency} อัตราส่วนไจโรแมกเนติกของโปรตอนคือ $\gamma/2\pi = 42.577~\text{MHz/T}$
\begin{equation*}
\nu_L = \left( \frac{\gamma}{2\pi} \right) B_0 = (42.577~\text{MHz/T}) \times 3.0~\text{T} = 127.731~\text{MHz}
\end{equation*}

\textbf{ขั้นตอนที่ 2} คำนวณความกว้างของแถบความถี่ $\Delta \nu_L$
\begin{equation*}
\Delta \nu_L = \left( \frac{\gamma}{2\pi} \right) \Delta B = (42.577 \times 10^6~\text{Hz/T}) \times (1.0 \times 10^{-4}~\text{T}) = 4257.7~\text{Hz} \approx 4.26~\text{kHz}
\end{equation*}

\textbf{Check (ตรวจสอบ)} ---- ความถี่ $127.7\text{ MHz}$ อยู่ในช่วงคลื่นวิทยุความถี่สูงมาก (Very High Frequency, VHF) ซึ่งตรงกับย่านความถี่วิทยุ FM เชิงพาณิชย์ และคลื่นนี้มีความปลอดภัยสูงเนื่องจากเป็นรังสีชนิดไม่แตกตัวเป็นไอออน (Non-ionizing Radiation)

\textbf{Think (คิดต่อ)} ---- ในการสร้างภาพ MRI วิศวกรจะจงใจใส่สนามแม่เหล็กเกรเดียนต์เชิงพื้นที่ ($G_x, G_y, G_z$) ซ้อนทับเข้าไป ทำให้ความถี่ลาร์มอร์ขึ้นกับตำแหน่งพิกัดสามมิติโดยตรง $\omega(x, y, z) = \gamma (B_0 + \mathbf{G}\cdot\mathbf{r})$ การแปลงฟูเรียร์ของสัญญาณวิทยุที่แผ่ออกมาจึงสามารถถอดรหัสตำแหน่งและสร้างภาพตัดขวางอวัยวะภายในได้อย่างแม่นยำ
\end{psctbox}

% ==========================================================
% สคริปต์การคำนวณเชิงตัวเลขภาษา Python
% ==========================================================
\section{สคริปต์การคำนวณเชิงตัวเลขภาษา Python จำลองออร์บิทัลสามมิติ}
\index{การคำนวณเชิงตัวเลขด้วยไพทอน (Python Computation)}

เพื่อศึกษาการกระจายตัวของความน่าจะเป็นเชิงรัศมีและการจำลองออร์บิทัลของอะตอมไฮโดรเจนเชิงตัวเลข นิสิตสามารถรันโค้ดภาษา Python ด้านล่างนี้ ซึ่งใช้โมดูล `scipy.special` และ `matplotlib` ในการคำนวณฟังก์ชันพหุนามลาแกร์คู่ควบ ฮาร์มอนิกส์ทรงกลม และการอินทิเกรตเชิงตัวเลข

\begin{tcolorbox}[colback=slateGrey!5, colframe=darkNavy, title={\textbf{Python 3 Script ---- Hydrogen Orbitals \& Radial Density Solver}}]
\begin{lstlisting}[language=Python]
import numpy as np
import matplotlib.pyplot as plt
from scipy.special import genlaguerre, sph_harm
from scipy.integrate import quad

# ค่าคงที่ทางฟิสิกส์ (หน่วยอะตอม a0 = 1)
a0 = 1.0

def radial_wavefunction(n, l, r):
    """คำนวณฟังก์ชันคลื่นเชิงรัศมี R_nl(r) ของอะตอมไฮโดรเจน"""
    rho = 2.0 * r / (n * a0)
    norm = np.sqrt((2.0 / (n * a0))**3 * np.math.factorial(n - l - 1) / 
                   (2.0 * n * (np.math.factorial(n + l))**3))
    # พหุนามลาแกร์คู่ควบ L_{n-l-1}^{2l+1}(rho)
    laguerre = genlaguerre(n - l - 1, 2 * l + 1)
    return norm * np.exp(-rho / 2.0) * (rho**l) * laguerre(rho)

def radial_probability_density(n, l, r):
    """คำนวณความหนาแน่นความน่าจะเป็นเชิงรัศมี P(r) = r^2 |R_nl(r)|^2"""
    R = radial_wavefunction(n, l, r)
    return (r**2) * (R**2)

# พล็อตเปรียบเทียบ P(r) ของสถานะ 1s, 2s, 2p, 3s, 3p, 3d
r_vals = np.linspace(0, 25 * a0, 1000)
states = [(1, 0, '1s'), (2, 0, '2s'), (2, 1, '2p'), 
          (3, 0, '3s'), (3, 1, '3p'), (3, 2, '3d')]

plt.figure(figsize=(10, 6))
for n, l, label in states:
    p_vals = [radial_probability_density(n, l, r) for r in r_vals]
    plt.plot(r_vals, p_vals, label=f'State {label} (n={n}, l={l})', lw=1.8)

plt.title('Hydrogen Radial Probability Density P(r) = r^2 |R_{nl}(r)|^2')
plt.xlabel('Radius r / a_0')
plt.ylabel('Probability Density P(r)')
plt.grid(True, linestyle='--', alpha=0.6)
plt.legend()
plt.tight_layout()
plt.savefig('hydrogen_radial_density.png', dpi=300)
print('กราฟความหนาแน่นความน่าจะเป็นถูกบันทึกที่ hydrogen_radial_density.png')

# คำนวณค่าคาดหวังเชิงตัวเลข <r> และเปรียบเทียบกับสูตรทฤษฎี
print('\nเปรียบเทียบค่าคาดหวัง <r> เชิงตัวเลขเทียบกับสูตรวิเคราะห์ ----')
for n, l, label in states:
    # อินทิเกรตเชิงตัวเลข quad( r * P(r), 0, inf )
    num_exp_r, _ = quad(lambda r: r * radial_probability_density(n, l, r), 0, 100 * a0)
    # สูตรแม่นตรงตามทฤษฎี
    exact_exp_r = (a0 / 2.0) * (3 * n**2 - l * (l + 1))
    print(f'State {label}: Numerical = {num_exp_r:.4f} a0 | Exact = {exact_exp_r:.4f} a0')
\end{lstlisting}
\end{tcolorbox}

% ==========================================================
% สรุปสาระสำคัญประจำบทที่ 6
% ==========================================================
\section{สรุปสาระสำคัญประจำบทที่ 6}

\begin{itemize}
    \item \textbf{การแยกตัวแปรในพิกัดทรงกลม} ---- สำหรับศักย์สู่ศูนย์กลาง $V(r)$ สมการชเรอดิงเงอร์สามมิติสามารถแยกตัวแปรเชิงรัศมีและเชิงมุมได้อย่างเด็ดขาด $\psi(r, \theta, \phi) = R(r)Y(\theta, \phi)$ โดยมีค่าคงตัวของการแยกตัวแปรคือ $\ell(\ell+1)$
    \item \textbf{ฮาร์มอนิกส์ทรงกลมและโมเมนตัมเชิงมุม} ---- ฟังก์ชัน $Y_\ell^m(\theta, \phi)$ เป็นสถานะเจาะจงร่วมของตัวดำเนินการ $\hat{\mathbf{L}}^2$ และ $\hat{L}_z$ โดยมีค่าเจาะจงคือ $\ell(\ell+1)\hbar^2$ และ $m\hbar$ ตามลำดับ ฮาร์มอนิกส์ทรงกลมกำหนดรูปร่างและสมมาตรเชิงมุมของออร์บิทัล
    \item \textbf{ระดับพลังงานของอะตอมไฮโดรเจน} ---- การบังคับให้ฟังก์ชันคลื่นลู่เข้าที่ระยะอนันต์นำไปสู่เลขควอนตัมหลัก $n \in \{1, 2, 3, \dots\}$ และสูตรพลังงาน $E_n = -13.606 Z^2/n^2~\text{eV}$ โดยมีความเสื่อมเชิงพื้นที่เท่ากับ $n^2$ ซึ่งเกิดจากสมมาตรซ่อนเร้น $SO(4)$ ของเวกเตอร์รุนเก-เลนซ์
    \item \textbf{ความหนาแน่นความน่าจะเป็นในแนวรัศมี} ---- ฟังก์ชัน $P(r) = r^2 |R_{n\ell}(r)|^2$ บ่งชี้โอกาสในการค้นพบอิเล็กตรอนในเปลือกทรงกลมหนา $dr$ โดยมีจำนวนบัพเชิงรัศมีเท่ากับ $n - \ell - 1$
    \item \textbf{การทดลองชแตร์น-แกร์ลัค} ---- ยืนยันการมีอยู่จริงของสปินอิเล็กตรอน $s = 1/2$ และการควอนไตซ์เชิงพื้นที่ของขั้วคู่แม่เหล็กในสนามแม่เหล็กไม่สม่ำเสมอ ซึ่งแยกออกเป็นสองสถานะ $m_s = \pm 1/2$
    \item \textbf{พีชคณิตสปินเพาลี} ---- ตัวดำเนินการสปินสัมพันธ์กับเมทริกซ์เพาลี $\hat{\mathbf{S}} = \frac{\hbar}{2}\boldsymbol{\sigma}$ ซึ่งมีสมบัติความเป็นแฮร์มิเชียน เอกภาพ และการแอนติคอมมิวต์
    \item \textbf{คาบการหมุน $4\pi$ ของสปินอร์} ---- การหมุนอนุภาคสปินครึ่งในปริภูมิสามมิติครบ $360^\circ$ จะทำให้สปินอร์เปลี่ยนเครื่องหมายเป็นลบ และต้องหมุนครบ $720^\circ$ จึงจะกลับคืนสู่สถานะเดิม
    \item \textbf{โครงสร้างละเอียด} ---- ผลรวมของการแก้ไขพลังงานจลน์เชิงสัมพัทธภาพ อันตรกิริยาสปิน-ออร์บิท ($\hat{\mathbf{L}}\cdot\hat{\mathbf{S}}$) และเทอมดาร์วิน แยกความเสื่อมของระดับพลังงานตามเลขควอนตัมโมเมนตัมเชิงมุมรวม $j = \ell \pm 1/2$
    \item \textbf{การประยุกต์เทคโนโลยีสมัยใหม่} ---- ทฤษฎีสปินเป็นรากฐานของเทคโนโลยีสปินทรอนิกส์ หัวอ่านความต้านทานแม่เหล็กยักษ์ GMR หน่วยความจำ MRAM ยุคใหม่ และเครื่องตรวจภาพเรโซแนนซ์แม่เหล็ก MRI ทางการแพทย์และเกษตรแม่นยำ
\end{itemize}

% ==========================================================
% แบบฝึกหัดท้ายบทที่ 6
% ==========================================================
\section*{แบบฝึกหัดท้ายบทที่ 6}
\addcontentsline{toc}{section}{แบบฝึกหัดท้ายบทที่ 6}

\subsection*{คำถามเชิงแนวคิด (Conceptual Questions)}

\begin{enumerate}[label=\textbf{\arabic*.}]
    \item เหตุใดสมการชเรอดิงเงอร์สำหรับศักย์สู่ศูนย์กลาง $V(r)$ จึงสามารถแยกตัวแปรส่วนมุม $(\theta, \phi)$ ออกจากตัวแปรระยะทาง $r$ ได้เสมอ ไม่ว่าฟังก์ชัน $V(r)$ จะมีรูปแบบใด
    \item อธิบายความหมายทางกายภาพของกำแพงหนีศูนย์กลางควอนตัม (Centrifugal Barrier) ในศักย์ยังผล $V_{\text{eff}}(r)$ เหตุใดอิเล็กตรอนที่มี $\ell > 0$ จึงไม่สามารถเข้าใกล้จุดศูนย์กลางของนิวเคลียสได้
    \item จงอธิบายความแตกต่างระหว่างความเสื่อมของระดับพลังงานตามเลขควอนตัม $m$ กับความเสื่อมตามเลขควอนตัม $\ell$ ในอะตอมไฮโดรเจน เหตุใดความเสื่อมตาม $\ell$ จึงถูกเรียกว่าความเสื่อมโดยบังเอิญ
    \item สำหรับสถานะ $1s$ ของอะตอมไฮโดรเจน ตำแหน่งที่มีความหนาแน่นความน่าจะเป็นต่อหน่วยปริมาตรสูงสุด ($|\psi|^2$ สูงสุด) กับตำแหน่งที่มีความหนาแน่นความน่าจะเป็นในแนวรัศมีสูงสุด ($P(r)$ สูงสุด) อยู่ที่ตำแหน่งเดียวกันหรือไม่ จงอธิบาย
    \item เหตุใดการทดลองชแตร์น-แกร์ลัคดั้งเดิมจึงต้องใช้ไอของอะตอมเงิน แทนที่จะใช้ลำอิเล็กตรอนอิสระในการยิงผ่านสนามแม่เหล็กไม่สม่ำเสมอ
    \item ในการจัดวางเครื่องมือชแตร์น-แกร์ลัคแบบเรียงลำดับ SGz $\to$ SGx $\to$ SGz เหตุใดลำอนุภาคจึงเกิดการแยกตัวออกเป็นสองแถบอีกครั้งในเครื่องที่สาม ทั้งที่ลำอนุภาคเคยถูกคัดเลือกให้มีสปินชี้ขึ้นตามแกน $z$ มาก่อนแล้ว
    \item จงอธิบายความหมายทางกายภาพของสปินอร์ และเหตุใดการหมุนสปินอร์ไปหนึ่งรอบเต็ม $360^\circ$ ในโลกกายภาพ จึงทำให้สปินอร์มีเครื่องหมายติดลบ
    \item เปรียบเทียบความแตกต่างระหว่างโมเมนตัมเชิงมุมออร์บิทัล $\hat{\mathbf{L}}$ กับโมเมนตัมเชิงมุมสปิน $\hat{\mathbf{S}}$ ในแง่ของปริภูมิการนิยามและค่าของเลขควอนตัมที่เป็นไปได้
    \item เหตุใดตัวดำเนินการ $\hat{L}_z$ และ $\hat{S}_z$ จึงไม่สลับที่กับฮามิลโทเนียนที่มีเทอมอันตรกิริยาสปิน-ออร์บิท และเหตุใดโมเมนตัมเชิงมุมรวม $\hat{\mathbf{J}} = \hat{\mathbf{L}} + \hat{\mathbf{S}}$ จึงเป็นตัวดำเนินการที่สลับที่
    \item อธิบายกลไกทางกายภาพของปรากฏการณ์ความต้านทานแม่เหล็กยักษ์ (GMR) และชี้ให้เห็นว่าสปินของอิเล็กตรอนควบคุมสภาพความต้านทานไฟฟ้าได้อย่างไร
    \item ในการตรวจวินิจฉัยทางการแพทย์ด้วย MRI เหตุใดจึงเลือกใช้นิวเคลียสของอะตอมไฮโดรเจนเป็นหลักในการสร้างภาพเนื้อเยื่อ
    \item ความถี่ลาร์มอร์ (Larmor Frequency) มีความสัมพันธ์กับขนาดของสนามแม่เหล็กภายนอกอย่างไร และวิศวกรใช้สนามแม่เหล็กเกรเดียนต์ในการระบุพิกัดสามมิติของภาพได้อย่างไร
    \item อธิบายความแตกต่างระหว่างเวลาผ่อนคลาย $T_1$ (Spin-Lattice) และ $T_2$ (Spin-Spin) ในปรากฏการณ์เรโซแนนซ์แม่เหล็กนิวเคลียร์
    \item เหตุใดค่าคาดหวังของระยะทาง $\langle r \rangle_{1s} = 1.5 a_0$ จึงมีค่ามากกว่ารัศมีที่มีโอกาสพบอิเล็กตรอนมากที่สุด $r_{\text{max}} = a_0$
    \item หากค่าคงตัวของพลังค์มีค่าเข้าใกล้ศูนย์ ($\hbar \to 0$) ผลการทดลองชแตร์น-แกร์ลัคจะเปลี่ยนแปลงไปอย่างไร
    \item อธิบายบทบาทของหลักพาริตี (Parity) ในฟังก์ชันฮาร์มอนิกส์ทรงกลม และการประยุกต์ใช้ในกฎการเลือกของการแผ่รังสีเชิงแสง
    \item เทอมดาร์วิน (Darwin Term) ในโครงสร้างละเอียดมีผลเฉพาะกับสถานะใด และมีคำอธิบายทางกายภาพจากการสั่นแบบซิทเทอร์เบเวกุง (Zitterbewegung) อย่างไร
    \item ทำไมหน่วยความจำ MRAM ที่ใช้สปินของอิเล็กตรอนจึงประหยัดพลังงานและทำงานเร็วกว่าหน่วยความจำแฟลชดั้งเดิม
    \item จงอธิบายว่าทำไมเวกเตอร์ลาปลาซ-รุนเก-เลนซ์ จึงเป็นปริมาณอนุรักษ์เฉพาะในศักย์ที่แปรผกผันกับระยะทางกำลังหนึ่ง ($1/r$) เท่านั้น
    \item ในการเกษตรแม่นยำยุคใหม่ เทคโนโลยี Portable NMR อาศัยคุณสมบัติสปินของนิวเคลียสในการตรวจวัดความชื้นในดินและพืชผลได้อย่างไร
\end{enumerate}

\subsection*{โจทย์ปัญหาเชิงวิเคราะห์และคำนวณ (Analytical Problems)}

\begin{enumerate}[label=\textbf{\arabic*.}]
    \item จงแสดงว่าฟังก์ชันฮาร์มอนิกส์ทรงกลม $Y_1^0(\theta, \phi) = \sqrt{\frac{3}{4\pi}}\cos\theta$ สอดคล้องกับเงื่อนไขการนอร์มอลไลซ์บนผิวทรงกลมสองมิติ $\int_0^{2\pi} d\phi \int_0^\pi d\theta \sin\theta |Y_1^0|^2 = 1$
    \item จงพิสูจน์โดยตรงในพิกัดคาร์ทีเซียนว่า $[\hat{L}_y, \hat{L}_z] = i\hbar\hat{L}_x$ โดยใช้คอมมิวเตเตอร์คาโนนิคัล $[\hat{x}_i, \hat{p}_j] = i\hbar\delta_{ij}$
    \item กำหนดตัวดำเนินการลดระดับโมเมนตัมเชิงมุม $\hat{L}_- = \hat{L}_x - i\hat{L}_y$ จงคำนวณผลลัพธ์ของ $[\hat{L}_z, \hat{L}_-]$ และ $[\hat{\mathbf{L}}^2, \hat{L}_-]$
    \item จงหาตำแหน่งของจุดบัพเชิงรัศมีทั้งหมดของออร์บิทัล $3s$ ของอะตอมไฮโดรเจน จากฟังก์ชัน $R_{30}(r)$ ในตารางที่~\ref{tab:radial_wavefunctions}
    \item สำหรับสถานะ $2p$ ($n=2, \ell=1$) ของอะตอมไฮโดรเจน จงหาระยะรัศมี $r_{\text{max}}$ ที่มีโอกาสพบอิเล็กตรอนมากที่สุด และเปรียบเทียบกับรัศมีวงโคจรของบอร์สำหรับ $n=2$ ($r_{\text{Bohr}} = 4a_0$)
    \item จงคำนวณค่าคาดหวัง $\langle 1/r \rangle$ สำหรับสถานะ $1s$ ของอะตอมไฮโดรเจน และแสดงว่าพลังงานศักย์เฉลี่ยสอดคล้องกับทฤษฎีบทไวเรียล $\langle V \rangle = 2E_1$
    \item จงคำนวณค่าคาดหวังของกำลังสองของระยะทาง $\langle r^2 \rangle$ สำหรับสถานะพื้น $1s$ ของอะตอมไฮโดรเจน และหาค่าความไม่แน่นอนของระยะรัศมี $\Delta r = \sqrt{\langle r^2 \rangle - \langle r \rangle^2}$
    \item สำหรับไฮโดรเจนิกไอออนฮีเลียมไอออน ($\text{He}^+, Z=2$) ในสถานะพื้น จงหารัศมีที่มีโอกาสพบอิเล็กตรอนมากที่สุด และคำนวณพลังงานที่ต้องใช้ในการทำให้อิเล็กตรอนหลุดเป็นอิสระ (Ionization Energy)
    \item ในการทดลองชแตร์น-แกร์ลัค ลำอะตอมเงินมีความเร็วเฉลี่ย $v = 450~\text{m/s}$ เคลื่อนที่ผ่านแม่เหล็กยาว $L = 10~\text{cm}$ โดยมีเกรเดียนต์ของสนามแม่เหล็ก $\partial B_z/\partial z = 1.2 \times 10^3~\text{T/m}$ หากมวลของอะตอมเงินคือ $1.79 \times 10^{-25}~\text{kg}$ จงหาระยะการแยกตัวระหว่างสองแถบ ($\Delta z$) บนฉากรับที่วางอยู่ที่ปลายขั้วแม่เหล็กพอดี
    \item จงหาค่าเจาะจงและสปินอร์เจาะจงที่นอร์มอลไลซ์แล้วของเมทริกซ์เพาลี $\sigma_y = \begin{pmatrix} 0 & -i \\ i & 0 \end{pmatrix}$ พร้อมทั้งแสดงว่าสปินอร์เจาะจงทั้งสองตั้งฉากซึ่งกันและกัน
    \item กำหนดสปินอร์สถานะ $\chi = \frac{1}{\sqrt{13}} \begin{pmatrix} 2 \\ 3i \end{pmatrix}$
    \begin{enumerate}
        \item จงตรวจสอบว่าสถานะนี้ได้รับการนอร์มอลไลซ์แล้วหรือไม่
        \item จงคำนวณความน่าจะเป็นที่จะวัดสปินตามแนวแกน $z$ แล้วได้สปินขึ้น และสปินลง
        \item จงคำนวณค่าคาดหวัง $\langle S_x \rangle, \langle S_y \rangle, \langle S_z \rangle$
    \end{enumerate}
    \item จงพิสูจน์เอกลักษณ์เมทริกซ์เพาลี $(\boldsymbol{\sigma}\cdot\mathbf{a})(\boldsymbol{\sigma}\cdot\mathbf{b}) = (\mathbf{a}\cdot\mathbf{b}) I + i \boldsymbol{\sigma}\cdot(\mathbf{a}\times\mathbf{b})$ โดยใช้สมบัติการสลับที่และการแอนติคอมมิวต์
    \item อนุภาคสปินครึ่งตัวหนึ่งอยู่ในสนามแม่เหล็กสม่ำเสมอ $\mathbf{B} = B_0 \hat{\mathbf{k}}$ ฮามิลโทเนียนคือ $\hat{H} = -\gamma B_0 \hat{S}_z$ หากที่เวลา $t = 0$ อนุภาคอยู่ในสถานะสปินชี้ตามแกน $+x$ จงหาสถานะสปินอร์ $\chi(t)$ ที่เวลา $t$ ใดๆ และคำนวณ $\langle S_x \rangle(t)$
    \item จงคำนวณผลต่างพลังงานโครงสร้างละเอียดระหว่างระดับ $2P_{3/2}$ และ $2P_{1/2}$ ในอะตอมไฮโดรเจน และหาความถี่และความยาวคลื่นของโฟตอนที่ปล่อยออกมาจากการเปลี่ยนระดับพลังงานนี้
    \item ในระบบสองระดับสปิน ตัวดำเนินการเลื่อนระดับสปินนิยามโดย $\hat{S}_+ = \hat{S}_x + i\hat{S}_y$ และ $\hat{S}_- = \hat{S}_x - i\hat{S}_y$ จงเขียนเมทริกซ์ของ $\hat{S}_+$ และ $\hat{S}_-$ ในฐานสปินแกน $z$ และคำนวณ $[\hat{S}_+, \hat{S}_-]$
    \item อะตอมไฮโดรเจนถูกวางในสนามแม่เหล็กภายนอกขนาดอ่อน $B_{\text{ext}} = 0.5~\text{T}$ จงคำนวณการแยกตัวของระดับพลังงานซีมานของสถานะ $1s_{1/2}$
    \item สำหรับสถานะ $3d$ ($n=3, \ell=2$) ของอะตอมไฮโดรเจน จงระบุค่าโมเมนตัมเชิงมุมรวม $j$ ที่เป็นไปได้ทั้งหมด และคำนวณค่าคาดหวัง $\langle \hat{\mathbf{L}}\cdot\hat{\mathbf{S}} \rangle$ สำหรับแต่ละสถานะ
    \item จงคำนวณความถี่ลาร์มอร์ของนิวเคลียสคาร์บอน-13 ($^{13}\text{C}$) ในสนามแม่เหล็ก $B_0 = 7.0~\text{Tesla}$ กำหนดให้อัตราส่วนไจโรแมกเนติกของ $^{13}\text{C}$ คือ $\gamma/2\pi = 10.705~\text{MHz/T}$
    \item สำหรับออร์บิทัล $2s$ ของอะตอมไฮโดรเจน จงคำนวณความน่าจะเป็นที่อิเล็กตรอนจะอยู่ภายในบริเวณบัพเชิงรัศมี ($r \le 2a_0$)
    \item \textbf{โจทย์ประยุกต์การคำนวณทางฟิสิกส์คอมพิวเตอร์ (Computational Task)} ---- จงพัฒนาต่อยอดสคริปต์ Python ในบทเรียนเพื่อคำนวณการกระจายตัวของความน่าจะเป็นในสามมิติ $|\psi_{210}(r, \theta, \phi)|^2$ ในระนาบ $xz$ จากนั้นพล็อตกราฟเส้นชั้นความสูง (Contour Plot) แสดงรูปร่างทรงดัมเบลล์ของออร์บิทัล $2p_z$ พร้อมทั้งคำนวณค่าอินทิกรัลสามมิติเพื่อตรวจสอบการนอร์มอลไลซ์ของฟังก์ชันคลื่นรวม
\end{enumerate}
